% figures06.tex -- Chapter 6 figures (Congruence of Angles and Triangles).
%
% Built 2026-08-17 from sources/Geometry.pdf pp. 98-129 (book pp. 84-115) with
% structure checked against the iPad scans (scan06 p11-p24, scan07 p1-p24,
% scan08 p1-p10; see sources/scan-index/).
%
% Constrained points are DERIVED, never eyeballed:
%   parallels  -> equal ratios along the two sides   ($(X)!r!(Y)$)
%   altitudes / perpendicular feet                   ($(X)!(P)!(Y)$)
%   crossings  -> (intersection of A--B and C--D)
%   midpoints  -> ($(X)!0.5!(Y)$)
% so the hypothesis the book states in words holds exactly in the drawing.
% The same hypotheses are re-asserted numerically in tools/constraints/ch06.py.
%
% Macro names are chapter-prefixed English number words: \FIGVI<NUMBERWORD>.
% A macro spanning several consecutive figures concatenates the words
% (e.g. \FIGVITWOTHREEFOUR holds Figs 6-2, 6-3 and 6-4).
%
% Line treatment follows the book: solid where the book prints solid (given),
% dashed where the book prints dashed (the constructed D' in the congruence
% proofs, and the "moved" angle B' in the inequality figures).  This chapter is
% otherwise plain `given' line work -- the book draws these diagrams uniformly.

\usetikzlibrary{decorations.pathreplacing,arrows.meta,angles,quotes}

% Primed labels: a math prime sits high enough that the collision audit reads it
% as a separate word overlapping its own letter; a small drop keeps the pair in
% one word and is visually indistinguishable.  (Same fix as ch09.)
\newcommand{\VIpr}{\raisebox{-1.3pt}{\ensuremath{{}^{\prime}}}}
\newcommand{\VIppr}{\VIpr\VIpr}
\newcommand{\VIpppr}{\VIpr\VIpr\VIpr}

% Chapter-local label clearance.  A labelled point lying ON a line needs a
% little more room than brumfiel.sty's default or the line clips the box.
\tikzset{
  vlab/.append style={outer sep=2.9pt},
  slab/.append style={outer sep=2.9pt},
  klab/.append style={outer sep=2.9pt},
}

% ---------------------------------------------------------------- house marks
% The congruence apparatus this chapter runs on, defined once.
%
% All helpers are chapter-prefixed: \aa and friends collide with LaTeX
% built-ins (\aa is the Scandinavian ring-a), so every one carries the VI code.
%
% \VItk{A}{B}{n} -- n tick marks across the middle of segment AB, perpendicular
%                 to it.  The book uses one tick for one pair of congruent
%                 segments, two for the second pair.
\newcommand{\VItk}[3]{%
  \foreach \tki in {1,...,#3}{%
    \pgfmathsetmacro{\tkd}{(\tki-(#3+1)/2)*2.0}%
    \coordinate (tkm) at ($(#1)!0.5!(#2)$);
    \coordinate (tkt) at ($(tkm)!\tkd pt!(#2)$);
    \draw[tick] ($(tkt)!2.6pt!90:(#2)$) -- ($(tkt)!2.6pt!270:(#2)$);}}
%
% The angle marks are struck by hand rather than through the `angles' pic
% library, because that library always sweeps COUNTERCLOCKWISE from its first
% ray to its second.  Whenever a call happened to name the two sides the other
% way round, the pic drew the REFLEX angle -- a 331-degree sweep that renders
% as a near-complete circle swallowing the whole figure (this is what put the
% spurious circle in Fig. 6-13, and it silently mis-marked several others).
% Striking the arc from the two bearings and normalising the sweep into
% (-180,180] makes every call draw the minor angle no matter which side is
% named first, so the mark can never depend on argument order again.
%
% \VIarc{P}{V}{Q}{r} -- single arc of radius r at vertex V, from side VP to VQ.
\newcommand{\VIarc}[4]{%
  \begingroup
  \pgfmathanglebetweenpoints{\pgfpointanchor{#2}{center}}{\pgfpointanchor{#1}{center}}%
  \edef\VIaP{\pgfmathresult}%
  \pgfmathanglebetweenpoints{\pgfpointanchor{#2}{center}}{\pgfpointanchor{#3}{center}}%
  \edef\VIaQ{\pgfmathresult}%
  \pgfmathparse{\VIaQ-\VIaP-360*round((\VIaQ-\VIaP)/360)}%
  \edef\VIdA{\pgfmathresult}%
  \draw[fig, line width=0.55pt] ([shift=(\VIaP:#4)]#2)
       arc[start angle=\VIaP, end angle=\VIaP+\VIdA, radius=#4];
  \endgroup}
% \VIarcb{P}{V}{Q}{r} -- double arc (the book's second congruent-angle mark).
\newcommand{\VIarcb}[4]{%
  \VIarc{#1}{#2}{#3}{#4}%
  \VIarc{#1}{#2}{#3}{#4+2.1pt}}
% \VIarcvia{P}{V}{Q}{W}{r} -- arc at V from side VP to side VQ, sweeping the way
%   that passes through ray VW.  \VIarc always takes the MINOR angle, which is
%   right for a numbered angle but cannot draw the marks Fig. 6-94 needs: there
%   the book sweeps one arc from OA all the way round to OA' on the side B lies
%   on -- a straight angle in case 2 and a reflex one in case 3.  Naming the ray
%   the arc must pass through picks the branch without hard-coding bearings, so
%   the mark still follows the points if they move.
\newcommand{\VIarcvia}[5]{%
  \begingroup
  \pgfmathanglebetweenpoints{\pgfpointanchor{#2}{center}}{\pgfpointanchor{#1}{center}}%
  \edef\VIvP{\pgfmathresult}%
  \pgfmathanglebetweenpoints{\pgfpointanchor{#2}{center}}{\pgfpointanchor{#3}{center}}%
  \edef\VIvQ{\pgfmathresult}%
  \pgfmathanglebetweenpoints{\pgfpointanchor{#2}{center}}{\pgfpointanchor{#4}{center}}%
  \edef\VIvW{\pgfmathresult}%
  \pgfmathparse{mod(\VIvQ-\VIvP+720,360)}\edef\VIdQ{\pgfmathresult}%
  \pgfmathparse{mod(\VIvW-\VIvP+720,360)}\edef\VIdW{\pgfmathresult}%
  \pgfmathparse{\VIdW<\VIdQ ? \VIdQ : \VIdQ-360}\edef\VIsw{\pgfmathresult}%
  \draw[fig, line width=0.55pt] ([shift=(\VIvP:#5)]#2)
       arc[start angle=\VIvP, end angle=\VIvP+\VIsw, radius=#5];
  \endgroup}
% \VIrt{P}{V}{Q}{r} -- the small square marking a right angle.  Built from the
% two side directions directly, so it always lands inside the angle.
\newcommand{\VIrt}[4]{%
  \begingroup
  \coordinate (rt@p) at ($(#2)!#4!(#1)$);
  \coordinate (rt@q) at ($(#2)!#4!(#3)$);
  \draw[fig, line width=0.55pt] (rt@p) -- ($(rt@p)+(rt@q)-(#2)$) -- (rt@q);
  \endgroup}
% \VInum{V}{ang}{d}{txt} -- a small numeral label (the book's "1", "2", ...
%   inside an angle), placed d from vertex V along direction ang.
\newcommand{\VInum}[4]{\node[slab] at ($(#1)+({#2}:#3)$) {#4};}
%
% \VInumb{P}{V}{Q}{d}{txt} -- the same numeral, but placed on the BISECTOR of
%   angle P-V-Q at distance d.  Placing by bisector rather than by a
%   hand-picked compass angle is what keeps the numeral off both rays, which
%   is where most label/geometry collisions in this chapter came from.
\newcommand{\VInumb}[5]{%
  \begingroup
  \coordinate (nb@1) at ($(#2)!1cm!(#1)$);
  \coordinate (nb@2) at ($(#2)!1cm!(#3)$);
  \coordinate (nb@m) at ($(nb@1)!0.5!(nb@2)$);
  \node[slab] at ($(#2)!#4!(nb@m)$) {#5};
  \endgroup}

% ================================================================== Fig 6-1
% Book p.84.  The given angle A (a free "V"), and below it the line l' with A'
% on it: r' is the ray of l' to the right of A', s' the constructed ray.  The
% dashed stub to the left of A' carries the caption "Given side of l'".
\newcommand{\FIGVIONE}{\begin{bkfigure}[\linewidth]
\begin{tikzpicture}[scale=1.235]
  % given angle A, drawn free above
  \coordinate (A)  at (2.55,1.32);
  \draw[given] (A) -- (3.12,2.14);
  \draw[given] (A) -- (3.95,0.92);
  \node[vlab, above left] at (A) {$A$};
  % line l' with A' on it
  \coordinate (Ap) at (1.72,0.16);
  \draw[given] (Ap) -- (3.62,-0.22);           % ray r'
  \draw[aux]   (Ap) -- (0.62,0.38);            % given side of l', dashed stub
  \draw[given] (Ap) -- (2.62,1.24);            % ray s'
  \node[vlab, below left] at (Ap) {$A\VIpr$};
  \node[slab, right] at (3.62,-0.22) {$r\VIpr$};
  \node[slab] at (2.10,1.02) {$s\VIpr$};
  \node[slab, below]  at (0.62,0.38) {$l\VIpr$};
  % The book runs the caption straight through the ray s', with the ray
  % passing in the word space.  Setting it as two nodes either side of the ray
  % reproduces that exactly and leaves the ray unobscured.
  \node[slab, anchor=east] at (2.01,0.60) {Given side};
  \node[slab, anchor=west] at (2.17,0.60) {of $l\VIpr$};
\end{tikzpicture}
\figcap{6--1}
\end{bkfigure}}

% =========================================================== Figs 6-2, 6-3, 6-4
% Book p.85, Exercise Group 6-1.
%  6-2: angle A with sides r (right) and s (up-left); ray r' from A' below.
%  6-3: angle A with sides r and s; a detached ray r' ending at A'.
%  6-4: a straight angle A on a line, with A' on it and ray r'.
\newcommand{\FIGVITWOTHREEFOUR}{\begin{bkfigure}[\linewidth]
\begin{minipage}{0.46\linewidth}\centering
\begin{tikzpicture}[scale=1.008]
  \coordinate (A)  at (1.35,0.52);
  \coordinate (Ap) at (0.22,-0.62);
  \draw[given] (A) -- (2.62,0.44);            % r
  \draw[given] (A) -- (0.52,1.02);            % toward r'/s corner
  \draw[given] (0.52,1.02) -- (0.30,1.62);    % s
  \draw[given] (Ap) -- (0.52,1.02);
  \node[vlab, below] at (A) {$A$};
  \node[vlab, below left] at (Ap) {$A\VIpr$};
  \node[slab, right] at (2.62,0.44) {$r$};
  \node[slab, above] at (0.30,1.62) {$s$};
  \node[slab, left]  at (0.52,1.02) {$r\VIpr$};
\end{tikzpicture}
\figcap{6--2}
\end{minipage}\hfill
\begin{minipage}{0.46\linewidth}\centering
\begin{tikzpicture}[scale=1.008]
  \coordinate (A) at (1.05,1.05);
  \draw[given] (A) -- (0.18,0.05);            % s
  \draw[given] (A) -- (1.72,0.12);            % r
  \node[vlab, above] at (A) {$A$};
  \node[slab, below left]  at (0.18,0.05) {$s$};
  \node[slab, below right] at (1.72,0.12) {$r$};
  % detached ray r' ending at A'
  \draw[given] (2.18,1.00) -- (3.30,1.00);
  \node[slab, above] at (2.30,1.00) {$r\VIpr$};
  \node[vlab, right] at (3.30,1.00) {$A\VIpr$};
\end{tikzpicture}
\figcap{6--3}

\medskip
\begin{tikzpicture}[scale=1.008]
  \coordinate (A)  at (0.95,0.72);
  \coordinate (Ap) at (1.42,0.72);
  \draw[given] (0.12,0.28) -- (Ap);
  \draw[given] (Ap) -- (3.05,0.16);
  \dt{A}
  \node[vlab, above] at (A) {$A$};
  \node[vlab, above right] at (Ap) {$A\VIpr$};
  \node[slab, right] at (3.05,0.16) {$r\VIpr$};
\end{tikzpicture}
\figcap{6--4}
\end{minipage}
\end{bkfigure}}

% ================================================================== Fig 6-5
% Book p.86.  THREE rays from each vertex, not four.  The common side splits
% the whole angle A into the two parts B and C.  The book marks this with two
% concentric arcs struck from the vertex: a SMALL one spanning the whole
% opening (that is angle A) and a LARGE one, also spanning the whole opening,
% which the middle ray cuts into the two arcs belonging to B and C.  Each
% letter is written just BEYOND its own arc -- A outside the small arc, B and
% C outside the large one, in their respective sectors.  The common side is
% drawn a little longer than the other two, as in the book.
% Ray directions and arc radii measured off the plate (book p.86, left fan):
% sides at 146/94/52 degrees, large arc at 0.89 and small arc at 0.35 of the
% side length.
\newcommand{\FIGVIFIVE}{\begin{bkfigure}[\linewidth]
\begin{tikzpicture}[scale=1.300]
 \foreach \dx/\lA/\lB/\lC in {0/A/B/C, 4.27/A\VIpr/B\VIpr/C\VIpr}{
  \begin{scope}[xshift=\dx cm]
   \coordinate (V) at (0,0);
   \draw[given] (V) -- ++(146:2.30);        % outer side, left
   \draw[given] (V) -- ++(94:2.38);         % common side of B and C
   \draw[given] (V) -- ++(52:2.20);         % outer side, right
   % large arc -- the middle ray cuts it into the marks for B and C
   \draw[tick] ([shift=(146:2.04)]V)
        arc[start angle=146, end angle=52, radius=2.04];
   % small arc near the vertex -- the mark for the whole angle A
   \draw[tick] ([shift=(146:0.80)]V)
        arc[start angle=146, end angle=52, radius=0.80];
   \node[vlab] at ($(V)+(116:1.14)$) {$\lA$};
   \node[vlab] at ($(V)+(120:2.36)$) {$\lB$};
   \node[vlab] at ($(V)+(71:2.33)$)  {$\lC$};
  \end{scope}}
\end{tikzpicture}
\figcap{6--5}
\end{bkfigure}}

% =========================================================== Figs 6-6, 6-7
% Book p.86, Exercise Group 6-3.
%  6-6: same two fans as 6-5 but labelled A,B,C and A',C',B' (the primed fan
%       has C' inside B', which is what makes part (b) of the exercise work).
%  6-7: two separate three-ray fans, P Q R from O and R V U from S T U -- the
%       book draws them as two "V" clusters sharing a baseline.
% 6-6 is the same three-ray idea as 6-5, but here each of the three angles
% carries its OWN arc at its own radius: the whole angle A on the outside, and
% the two parts on smaller arcs.  The unprimed fan opens up-left to
% horizontal-right with B the upper part and C the lower; the primed fan is
% the mirror-ish case with C' the part next to the horizontal side and B' the
% part next to the slanting one -- that swap is what the exercise turns on.
\newcommand{\FIGVISIXSEVEN}{\begin{bkfigure}[\linewidth]
\begin{minipage}{0.60\linewidth}\centering
\begin{tikzpicture}[scale=1.040]
   % ---- unprimed fan: sides at 134 / 89 / 0 degrees
   \coordinate (V) at (0,0);
   \draw[given] (V) -- ++(134:1.84);
   \draw[given] (V) -- ++(89:1.87);
   \draw[given] (V) -- ++(0:1.86);
   \draw[tick] ([shift=(134:1.61)]V)
        arc[start angle=134, end angle=0, radius=1.61];    % whole angle A
   \draw[tick] ([shift=(134:0.50)]V)
        arc[start angle=134, end angle=89, radius=0.50];   % part B
   \draw[tick] ([shift=(89:1.06)]V)
        arc[start angle=89, end angle=0, radius=1.06];     % part C
   \node[vlab] at ($(V)+(64:1.94)$)  {$A$};
   \node[vlab] at ($(V)+(111:0.96)$) {$B$};
   \node[vlab] at ($(V)+(45:0.62)$)  {$C$};
   % ---- primed fan: sides at 163 / 89 / 48 degrees
   \coordinate (W) at (4.24,0);
   \draw[given] (W) -- ++(163:1.91);
   \draw[given] (W) -- ++(89:2.06);
   \draw[given] (W) -- ++(48:1.93);
   \draw[tick] ([shift=(163:1.58)]W)
        arc[start angle=163, end angle=48, radius=1.58];   % whole angle A'
   \draw[tick] ([shift=(163:1.06)]W)
        arc[start angle=163, end angle=89, radius=1.06];   % part C'
   \draw[tick] ([shift=(89:0.50)]W)
        arc[start angle=89, end angle=48, radius=0.50];    % part B'
   \node[vlab] at ($(W)+(119:2.01)$) {$A\VIpr$};
   \node[vlab] at ($(W)+(126:0.62)$) {$C\VIpr$};
   \node[vlab] at ($(W)+(68:1.10)$)  {$B\VIpr$};
\end{tikzpicture}
\figcap{6--6}
\end{minipage}\hfill
\begin{minipage}{0.36\linewidth}\centering
\begin{tikzpicture}[scale=0.832]
  \coordinate (O) at (0.72,0);
  \draw[given] (O) -- (0.10,1.62);   \node[vlab, above] at (0.10,1.62) {$P$};
  \draw[given] (O) -- (0.62,1.62);   \node[vlab, above] at (0.62,1.62) {$Q$};
  \draw[given] (O) -- (1.46,1.62);   \node[vlab, above] at (1.46,1.62) {$R$};
  \node[vlab, below] at (O) {$O$};
  \coordinate (S) at (1.72,0);
  \coordinate (T) at (2.10,0);
  \coordinate (U) at (2.72,0);
  \coordinate (Vv) at (2.16,1.62);
  \draw[given] (S) -- (Vv);
  \draw[given] (T) -- (Vv);
  \draw[given] (U) -- (Vv);
  \node[vlab, below] at (S) {$S$};
  \node[vlab, below] at (T) {$T$};
  \node[vlab, below] at (U) {$U$};
  \node[vlab, above] at (Vv) {$V$};
\end{tikzpicture}
\figcap{6--7}
\end{minipage}
\end{bkfigure}}

% =========================================================== Figs 6-8, 6-9
% Book p.87.  6-8: fan at A with rays to B, C, D, E.
%             6-9: fan at P with rays to R, S, T, U.
\newcommand{\FIGVIEIGHTNINE}{\begin{bkfigure}[\linewidth]
\begin{minipage}{0.48\linewidth}\centering
\begin{tikzpicture}[scale=1.086]
  \coordinate (A) at (1.30,0);
  \draw[given] (A) -- (0.10,1.28);  \node[vlab, above left]  at (0.10,1.28) {$B$};
  \draw[given] (A) -- (1.02,1.62);  \node[vlab, above] at (1.02,1.62) {$C$};
  \draw[given] (A) -- (1.42,1.62);  \node[vlab, above] at (1.42,1.62) {$D$};
  \draw[given] (A) -- (2.42,1.10);  \node[vlab, right] at (2.42,1.10) {$E$};
  \node[vlab, below] at (A) {$A$};
\end{tikzpicture}
\figcap{6--8}
\end{minipage}\hfill
\begin{minipage}{0.48\linewidth}\centering
\begin{tikzpicture}[scale=1.086]
  \coordinate (P) at (1.30,0);
  \draw[given] (P) -- (0.28,1.42);  \node[vlab, above left] at (0.28,1.42) {$R$};
  \draw[given] (P) -- (0.92,1.62);  \node[vlab, above] at (0.92,1.62) {$S$};
  \draw[given] (P) -- (1.52,1.62);  \node[vlab, above] at (1.52,1.62) {$T$};
  \draw[given] (P) -- (2.16,1.44);  \node[vlab, above right] at (2.16,1.44) {$U$};
  \node[vlab, below] at (P) {$P$};
\end{tikzpicture}
\figcap{6--9}
\end{minipage}
\end{bkfigure}}

% ==================================================== Figs 6-10, 6-11, 6-12
% Book p.87.  6-10: two supplement pictures (angle A with its adjacent
%   supplement B on a line; the primed copy).  6-11: two crossing lines with
%   vertical angles A and A', and B one of the others.  6-12: a right angle.
\newcommand{\FIGVITENELEVENTWELVE}{\begin{bkfigure}[\linewidth]
\begin{minipage}{0.50\linewidth}\centering
% In the book the ray rises to the RIGHT from a point on the line, so the
% OBTUSE angle (marked by the big arc sweeping over the top from the
% left-hand ray) is A and the small acute one tucked between the ray and the
% right-hand part of the line is B.  An earlier draft mirrored the ray and
% then put both letters in the same wedge, so the picture no longer showed a
% pair of supplements at all.
\begin{tikzpicture}
 \foreach \dx/\lA/\lB in {0/A/B, 3.30/A\VIpr/B\VIpr}{
  \begin{scope}[xshift=\dx cm]
   \coordinate (O)  at (0,0);
   \coordinate (Lt) at (-1.05,0);
   \coordinate (Rt) at (1.55,0);
   \coordinate (St) at (32:1.55);
   \draw[given] (Lt) -- (Rt);
   \draw[given] (O) -- (St);
   \VIarc{Lt}{O}{St}{0.62cm}   % the obtuse angle A
   \VIarc{St}{O}{Rt}{0.40cm}   % the acute angle B
   \node[slab] at ($(O)+(118:0.86)$) {$\lA$};
   \node[slab] at ($(O)+(16:1.02)$)  {$\lB$};
  \end{scope}}
\end{tikzpicture}
\figcap{6--10}
\end{minipage}\hfill
\begin{minipage}{0.30\linewidth}\centering
% Two crossing lines: A and A' are the vertical pair (left and right), B is
% one of the other two (the top one).  Each letter sits on the BISECTOR of
% its own wedge -- off the bisector a letter drifts into one of the four rays,
% which is what the collision audit was catching here.
\begin{tikzpicture}
  \coordinate (O)  at (0,0);
  \coordinate (P1) at (-1.30,-0.44);   \coordinate (P2) at (1.30,0.44);
  \coordinate (Q1) at (-1.30,0.44);    \coordinate (Q2) at (1.30,-0.44);
  \draw[given] (P1) -- (P2);
  \draw[given] (Q1) -- (Q2);
  \VIarc{Q1}{O}{P1}{0.60cm}   % left vertical angle  A
  \VIarc{P2}{O}{Q1}{0.90cm}   % top angle            B
  \VIarc{P2}{O}{Q2}{0.60cm}   % right vertical angle A'
  \node[slab] at ($(O)+(208:1.45)$) {$A$};
  \node[slab] at ($(O)+(90:0.70)$)  {$B$};
  \node[slab] at ($(O)+(-32:1.56)$)   {$A\VIpr$};
\end{tikzpicture}
\figcap{6--11}
\end{minipage}\hfill
\begin{minipage}{0.16\linewidth}\centering
\begin{tikzpicture}[scale=0.80]
  \coordinate (O) at (0.62,0);
  \coordinate (R) at (1.24,0);
  \coordinate (U) at (0.62,1.02);
  \draw[given] (0.02,0) -- (R);
  \draw[given] (O) -- (U);
  \VIrt{R}{O}{U}{6pt}
\end{tikzpicture}
\figcap{6--12}
\end{minipage}
\end{bkfigure}}

% ==================================================== Figs 6-13, 6-14, 6-15
% Book p.88, Exercise Group 6-4.
%  6-13: fan at A with rays to E, D, C, B; angle 1 = DAE, angle 2 = CAB.
%  6-14: two parallel lines cut by a transversal through A, B, C, D, E.
%  6-15: three concurrent lines through O: A--F, C--D, E--B.
\newcommand{\FIGVITHIRTEENFOURTEENFIFTEEN}{\begin{bkfigure}[\linewidth]
\begin{minipage}{0.46\linewidth}\centering
\begin{tikzpicture}[scale=1.009]
  \coordinate (A) at (1.42,0);
  \coordinate (E) at (0.05,1.02);
  \coordinate (D) at (0.72,1.52);
  \coordinate (C) at (2.02,1.52);
  \coordinate (B) at (2.62,1.02);
  \draw[given] (A) -- (E);  \draw[given] (A) -- (D);
  \draw[given] (A) -- (C);  \draw[given] (A) -- (B);
  \node[vlab, left]  at (E) {$E$};
  \node[vlab, above] at (D) {$D$};
  \node[vlab, above] at (C) {$C$};
  \node[vlab, right] at (B) {$B$};
  \node[vlab, below] at (A) {$A$};
  \VIarc{E}{A}{D}{0.86cm}
  \VIarc{C}{A}{B}{0.86cm}
  % Both wedges here are under 30 degrees wide, so a numeral set at a
  % hand-picked bearing lands on a ray; take the bisector instead.
  \VInumb{E}{A}{D}{1.10cm}{1}
  \VInumb{C}{A}{B}{1.10cm}{2}
\end{tikzpicture}
\figcap{6--13}

\medskip
% Rebuilt against scan06 p17 / photo p.88.  A, B, C, D, E are ANGLES, not
% points: the book draws THREE parallels cut by one transversal and strikes an
% arc on five of the angles so formed -- A below-left and B above-right at the
% lowest crossing, C below-right at the middle one, D below-right and E
% above-right at the top one.  (An earlier draft drew two parallels and
% lettered A..E as five POINTS, which left Exercise 3 -- "angle B is congruent
% to the supplement of angle C, and angle C is congruent to angle D; prove
% angle A is congruent to angle E" -- with nothing in the picture to name.)
\begin{tikzpicture}[scale=1.009]
  \coordinate (L1a) at (0.02,0.05);  \coordinate (L1b) at (3.20,0.05);
  \coordinate (L2a) at (0.02,1.00);  \coordinate (L2b) at (3.20,1.00);
  \coordinate (L3a) at (0.02,1.95);  \coordinate (L3b) at (3.20,1.95);
  \draw[given] (L1a) -- (L1b);
  \draw[given] (L2a) -- (L2b);
  \draw[given] (L3a) -- (L3b);
  \coordinate (P) at (1.02,-0.28);
  \coordinate (Q) at (2.86,2.24);
  \draw[given] (P) -- (Q);
  \coordinate (X1) at (intersection of P--Q and L1a--L1b);
  \coordinate (X2) at (intersection of P--Q and L2a--L2b);
  \coordinate (X3) at (intersection of P--Q and L3a--L3b);
  % arcs: the angle each letter names
  \VIarc{L1a}{X1}{P}{0.30cm}    % A -- below-left at the lowest crossing
  \VIarc{L1b}{X1}{Q}{0.34cm}    % B -- above-right there
  \VIarc{L2b}{X2}{P}{0.34cm}    % C -- below-right at the middle crossing
  \VIarc{L3b}{X3}{P}{0.34cm}    % D -- below-right at the top crossing
  \VIarc{L3b}{X3}{Q}{0.34cm}    % E -- above-right there
  \VInumb{L1a}{X1}{P}{0.56cm}{$A$}
  \VInumb{L1b}{X1}{Q}{0.60cm}{$B$}
  \VInumb{L2b}{X2}{P}{0.60cm}{$C$}
  \VInumb{L3b}{X3}{P}{0.60cm}{$D$}
  \VInumb{L3b}{X3}{Q}{0.62cm}{$E$}
\end{tikzpicture}
\figcap{6--14}
\end{minipage}\hfill
\begin{minipage}{0.46\linewidth}\centering
\begin{tikzpicture}[scale=1.009]
  % The six ends carry names so the constraint suite can read them straight
  % out of the drawing; the three lines really do meet at O.
  \coordinate (A) at (0.10,1.62);
  \coordinate (F) at (2.74,0.10);
  \coordinate (E) at (0.10,0.10);
  \coordinate (B) at (2.74,1.62);
  \coordinate (C) at (1.42,1.72);
  \coordinate (D) at (1.42,0.02);
  \coordinate (O) at (intersection of A--F and E--B);
  \draw[given] (A) -- (F);
  \draw[given] (E) -- (B);
  \draw[given] (C) -- (D);
  \node[vlab, above left]  at (A) {$A$};
  \node[vlab, above] at (C) {$C$};
  \node[vlab, above right] at (B) {$B$};
  \node[vlab, below left]  at (E) {$E$};
  \node[vlab, below] at (D) {$D$};
  \node[vlab, below right] at (F) {$F$};
  % Three lines meet at O, so six rays leave it; the letter goes on the
  % bisector of one of the 60-degree gaps rather than beside a ray.
  \node[vlab] at ($(O)+(68:0.64)$) {$O$};
\end{tikzpicture}
\figcap{6--15}
\end{minipage}
\end{bkfigure}}

% ================================================================= Fig 6-16
% Book p.89.  The S.A.S. postulate picture: triangles ABC and A'B'C' with one
% tick on AB/A'B', two on AC/A'C', and a double arc at A and A'.
\newcommand{\FIGVISIXTEEN}{\begin{bkfigure}[\linewidth]
\begin{tikzpicture}[scale=1.235]
 \foreach \dx/\lA/\lB/\lC in {0/A/B/C, 3.85/A\VIpr/B\VIpr/C\VIpr}{
  \begin{scope}[xshift=\dx cm]
   \coordinate (A) at (0,0);
   \coordinate (C) at (2.72,0);
   \coordinate (B) at (1.18,1.42);
   \draw[given] (A) -- (C) -- (B) -- cycle;
   \VItk{A}{B}{1}
   \VItk{A}{C}{2}
   \VIarcb{C}{A}{B}{0.52cm}
   \node[vlab, below left]  at (A) {$\lA$};
   \node[vlab, above] at (B) {$\lB$};
   \node[vlab, below right] at (C) {$\lC$};
  \end{scope}}
\end{tikzpicture}
\figcap{6--16}
\end{bkfigure}}

% =========================================================== Figs 6-17, 6-18
% Book p.89, Exercise Group 6-5.
%  6-17: triangles ABC (flat, opening right) and DEF.
%  6-18: triangles PQR and STU.
\newcommand{\FIGVISEVENTEENEIGHTEEN}{\begin{bkfigure}[\linewidth]
\begin{minipage}{0.50\linewidth}\centering
\begin{tikzpicture}[scale=0.894]
  \coordinate (A) at (0,0.62);
  \coordinate (B) at (1.02,0);
  \coordinate (C) at (1.92,0.68);
  \draw[given] (A) -- (B) -- (C) -- cycle;
  \node[vlab, left]  at (A) {$A$};
  \node[vlab, below] at (B) {$B$};
  \node[vlab, right] at (C) {$C$};
  % The two triangles are separate figures; keep a clear gap between C and E
  % so their letters cannot run together.
  \coordinate (E) at (2.62,0);
  \coordinate (D) at (3.28,1.42);
  \coordinate (F) at (4.24,1.10);
  \draw[given] (E) -- (D) -- (F) -- cycle;
  \node[vlab, below] at (E) {$E$};
  \node[vlab, above] at (D) {$D$};
  \node[vlab, right] at (F) {$F$};
\end{tikzpicture}
\figcap{6--17}
\end{minipage}\hfill
\begin{minipage}{0.46\linewidth}\centering
\begin{tikzpicture}[scale=0.894]
  \coordinate (P) at (0,0);
  \coordinate (R) at (1.32,0);
  \coordinate (Q) at (0.86,1.30);
  \draw[given] (P) -- (R) -- (Q) -- cycle;
  \node[vlab, left]  at (P) {$P$};
  \node[vlab, right] at (R) {$R$};
  \node[vlab, above] at (Q) {$Q$};
  % R and U were close enough that their letters merged into one word.
  \coordinate (U) at (2.14,0);
  \coordinate (S) at (2.04,1.30);
  \coordinate (T) at (3.38,1.30);
  \draw[given] (U) -- (S) -- (T) -- cycle;
  \node[vlab, below] at (U) {$U$};
  \node[vlab, above left]  at (S) {$S$};
  \node[vlab, above right] at (T) {$T$};
\end{tikzpicture}
\figcap{6--18}
\end{minipage}
\end{bkfigure}}

% =========================================================== Figs 6-19, 6-20
% Book p.89.  6-19: A above, B below it, with C and D at the ends -- the
%   "arrowhead" of Exercise 5 (AC = AD, angle CAB = angle BAD).
%  6-20: segments D--B and A--C crossing at O.
\newcommand{\FIGVININETEENTWENTY}{\begin{bkfigure}[\linewidth]
\begin{minipage}{0.46\linewidth}\centering
\begin{tikzpicture}[scale=0.966]
  \coordinate (A) at (1.42,1.52);
  \coordinate (B) at (1.42,0.52);
  \coordinate (C) at (0.02,0);
  \coordinate (D) at (2.82,0);
  \draw[given] (C) -- (A) -- (D);
  \draw[given] (C) -- (B) -- (D);
  \draw[given] (A) -- (B);
  \node[vlab, above] at (A) {$A$};
  \node[vlab, below] at (B) {$B$};
  \node[vlab, left]  at (C) {$C$};
  \node[vlab, right] at (D) {$D$};
\end{tikzpicture}
\figcap{6--19}
\end{minipage}\hfill
\begin{minipage}{0.50\linewidth}\centering
\begin{tikzpicture}[scale=0.966]
  \coordinate (D) at (0,0.52);
  \coordinate (B) at (3.02,0.42);
  \coordinate (A) at (0.86,0);
  \coordinate (C) at (2.42,1.02);
  \draw[given] (D) -- (B);
  \draw[given] (A) -- (C);
  \coordinate (O) at (intersection of D--B and A--C);
  \node[vlab, left]  at (D) {$D$};
  \node[vlab, right] at (B) {$B$};
  \node[vlab, below] at (A) {$A$};
  \node[vlab, above] at (C) {$C$};
  \node[vlab, below right] at (O) {$O$};
\end{tikzpicture}
\figcap{6--20}
\end{minipage}
\end{bkfigure}}

% =========================================================== Figs 6-21, 6-22
% Book p.90.  6-21: rectangle A B D C with the diagonal C--B.
%             6-22: triangle ACB with the cevian CD to the base.
\newcommand{\FIGVITWENTYONETWENTYTWO}{\begin{bkfigure}[\linewidth]
\begin{minipage}{0.50\linewidth}\centering
\begin{tikzpicture}[scale=0.991]
  \coordinate (A) at (0,0.92);
  \coordinate (B) at (2.92,0.92);
  \coordinate (C) at (0,0);
  \coordinate (D) at (2.92,0);
  \draw[given] (A) -- (B) -- (D) -- (C) -- cycle;
  \draw[given] (C) -- (B);
  \node[vlab, above left]  at (A) {$A$};
  \node[vlab, above right] at (B) {$B$};
  \node[vlab, below left]  at (C) {$C$};
  \node[vlab, below right] at (D) {$D$};
\end{tikzpicture}
\figcap{6--21}
\end{minipage}\hfill
\begin{minipage}{0.46\linewidth}\centering
\begin{tikzpicture}[scale=0.991]
  \coordinate (A) at (0,0);
  \coordinate (B) at (2.32,0);
  \coordinate (C) at (1.12,1.62);
  \coordinate (D) at ($(A)!0.48!(B)$);
  \draw[given] (A) -- (B) -- (C) -- cycle;
  \draw[given] (C) -- (D);
  \node[vlab, left]  at (A) {$A$};
  \node[vlab, right] at (B) {$B$};
  \node[vlab, above] at (C) {$C$};
  \node[vlab, below] at (D) {$D$};
\end{tikzpicture}
\figcap{6--22}
\end{minipage}
\end{bkfigure}}

% =========================================================== Figs 6-23, 6-24
% Book p.90.  6-23: A--B--C on a line with E and D above (two triangles meeting
%   at B).  6-24: quadrilateral A D C with B at the right, angles 1..4.
\newcommand{\FIGVITWENTYTHREETWENTYFOUR}{\begin{bkfigure}[\linewidth]
\begin{minipage}{0.52\linewidth}\centering
\begin{tikzpicture}[scale=0.940]
  \coordinate (A) at (0,0);
  \coordinate (B) at (1.62,0);
  \coordinate (C) at (3.12,0);
  \coordinate (E) at (0.86,1.22);
  \coordinate (D) at (2.02,1.28);
  \draw[given] (A) -- (C);
  \draw[given] (A) -- (E) -- (B);
  \draw[given] (B) -- (D) -- (C);
  \node[vlab, left]  at (A) {$A$};
  \node[vlab, below] at (B) {$B$};
  \node[vlab, right] at (C) {$C$};
  \node[vlab, above] at (E) {$E$};
  \node[vlab, above] at (D) {$D$};
\end{tikzpicture}
\figcap{6--23}
\end{minipage}\hfill
\begin{minipage}{0.44\linewidth}\centering
\begin{tikzpicture}[scale=1.490]
  \coordinate (A) at (0.72,1.52);
  \coordinate (C) at (0.72,0);
  % Exercise 10 gives AD = DC, so D must sit on the perpendicular bisector of
  % AC -- i.e. at the height of AC's midpoint.
  % Measured off scan06 p20 (book p.90): the book's quadrilateral is A D C B
  % with the single diagonal D--B.  There is NO segment A--C: D sits only a
  % little to the left of the line AC (offset 0.12 of AC's length), so the
  % angle ADC is nearly straight and its one big arc sweeps almost 180 deg.
  % An earlier draft drew A--C as well and pushed D out to 0.46 of AC, which
  % turned the plate into a diamond and buried the numerals 1 and 2 under the
  % spurious chord.
  \coordinate (D) at (0.54,0.76);
  \coordinate (B) at (2.62,0.76);
  \draw[given] (A) -- (D) -- (C);
  \draw[given] (A) -- (B) -- (C);
  \draw[given] (D) -- (B);
  \node[vlab, above] at (A) {$A$};
  \node[vlab, below] at (C) {$C$};
  \node[vlab, left]  at (D) {$D$};
  \node[vlab, right] at (B) {$B$};
  % D carries one big arc across the whole of angle ADC, which the segment DB
  % cuts into the parts 2 (upper) and 1 (lower); B carries a small arc for
  % each of 3 and 4.  Every numeral sits on its own bisector.
  \VIarc{A}{D}{C}{0.40cm}
  \VInumb{A}{D}{B}{0.62cm}{2}
  \VInumb{B}{D}{C}{0.62cm}{1}
  \VIarc{A}{B}{D}{0.52cm}   \VInumb{A}{B}{D}{1.20cm}{3}
  \VIarc{D}{B}{C}{0.52cm}   \VInumb{D}{B}{C}{1.20cm}{4}
\end{tikzpicture}
\figcap{6--24}
\end{minipage}
\end{bkfigure}}

% ==================================================== Figs 6-25, 6-26, 6-27
% Book p.90.
%  6-25: the zigzag B G C F H D A -- C and D are the upper vertices, G and H
%        the lower ones, with F between C and D on the upper line.
%  6-26: A B C on a line with peaks E and D; angles 1 and 2 at the ends.
%  6-27: triangle A M B with apex C, cevian CM, angles 1..4.
\newcommand{\FIGVITWENTYFIVETWENTYSIXTWENTYSEVEN}{\begin{bkfigure}[\linewidth]
\begin{tikzpicture}[scale=0.984]
  \coordinate (B) at (0,0);
  \coordinate (C) at (0.62,1.12);
  \coordinate (G) at (1.42,0);
  \coordinate (F) at (2.12,1.12);
  \coordinate (H) at (2.82,0);
  \coordinate (D) at (3.62,1.12);
  \coordinate (A) at (4.24,0);
  \draw[given] (B) -- (C) -- (G) -- (F) -- (H) -- (D) -- (A);
  \node[vlab, left]  at (B) {$B$};
  \node[vlab, above] at (C) {$C$};
  \node[vlab, below] at (G) {$G$};
  \node[vlab, above] at (F) {$F$};
  \node[vlab, below] at (H) {$H$};
  \node[vlab, above] at (D) {$D$};
  \node[vlab, right] at (A) {$A$};
  \VItk{C}{F}{2}
  \VItk{F}{D}{2}
  % Exercise 11 gives CG = HD as well; the book marks that pair with a single
  % tick, which an earlier draft left off.
  \VItk{C}{G}{1}
  \VItk{H}{D}{1}
  % Angle 3 is the one BETWEEN the top line and FG (i.e. angle CFG), not the
  % whole wedge GFH: the proof subtracts angle 1 from angle BCF to get angle
  % GCF, and it is the pair CFG / DFH that S.A.S. then makes congruent.
  \VIarc{B}{C}{G}{0.34cm}  \VInumb{B}{C}{G}{0.50cm}{1}
  \VIarc{C}{F}{G}{0.34cm}  \VInumb{C}{F}{G}{0.56cm}{3}
  \VIarc{H}{F}{D}{0.34cm}  \VInumb{H}{F}{D}{0.56cm}{4}
  \VIarc{H}{D}{A}{0.34cm}  \VInumb{H}{D}{A}{0.50cm}{2}
\end{tikzpicture}
\figcap{6--25}

\medskip
\begin{minipage}{0.52\linewidth}\centering
\begin{tikzpicture}[scale=0.938]
  % Exercise 12 gives AE = CD and AB = BC, so the two peaks are mirror
  % images about B and B is the midpoint of AC.  The draft had AB = 1.72
  % against BC = 1.70 and AE = 1.46 against CD = 1.60.
  \coordinate (A) at (0,0);
  \coordinate (C) at (3.42,0);
  \coordinate (B) at ($(A)!0.5!(C)$);
  \coordinate (E) at (0.62,1.32);
  \coordinate (D) at (2.80,1.32);
  \draw[given] (-0.28,0) -- (3.70,0);
  \draw[given] (A) -- (E) -- (B) -- (D) -- (C);
  \node[vlab] at ($(A)+(150:0.66)$) {$A$};
  \node[vlab, below] at (B) {$B$};
  \node[vlab] at ($(C)+(30:0.66)$) {$C$};
  \node[vlab, above] at (E) {$E$};
  \node[vlab, above] at (D) {$D$};
  % 1 and 2 are the EXTERIOR angles at the ends of the line -- between the
  % part of the line beyond A (resp. beyond C) and the side AE (resp. CD) --
  % and the book strikes an arc across each.  Those arcs were missing.
  \coordinate (Lx) at (-0.28,0);
  \coordinate (Rx) at (3.70,0);
  \VIarc{Lx}{A}{E}{0.30cm}   \VInum{A}{210}{0.52cm}{1}
  \VIarc{Rx}{C}{D}{0.30cm}   \VInum{C}{330}{0.52cm}{2}
\end{tikzpicture}
\figcap{6--26}
\end{minipage}\hfill
\begin{minipage}{0.44\linewidth}\centering
\begin{tikzpicture}[scale=0.938]
  \coordinate (A) at (0,0);
  \coordinate (B) at (2.62,0);
  \coordinate (C) at (1.32,1.72);
  \coordinate (M) at ($(A)!0.5!(B)$);
  % The book carries the base BEYOND A and beyond B, and angles 1 and 2 are
  % the EXTERIOR angles struck against those stubs -- not the interior base
  % angles an earlier draft marked (which, with AC = BC given, would have been
  % congruent by inspection and left Exercise 13 with nothing to prove).
  \coordinate (Lx) at (-0.52,0);
  \coordinate (Rx) at (3.14,0);
  \draw[given] (Lx) -- (Rx);
  \draw[given] (A) -- (C) -- (B);
  \draw[given] (C) -- (M);
  \node[vlab, below left]  at (A) {$A$};
  \node[vlab, below right] at (B) {$B$};
  \node[vlab, above] at (C) {$C$};
  \node[vlab, below] at (M) {$M$};
  \VIarc{A}{C}{M}{0.46cm}  \VInumb{A}{C}{M}{0.64cm}{3}
  \VIarc{M}{C}{B}{0.46cm}  \VInumb{M}{C}{B}{0.64cm}{4}
  \VIarc{Lx}{A}{C}{0.34cm} \VInumb{Lx}{A}{C}{0.60cm}{1}
  \VIarc{Rx}{B}{C}{0.34cm} \VInumb{Rx}{B}{C}{0.60cm}{2}
\end{tikzpicture}
\figcap{6--27}
\end{minipage}
\end{bkfigure}}

% =========================================================== Figs 6-28, 6-29
% Book p.91.  6-28: parallelogram-shaped quadrilateral D C B A with diagonal
%   D--B, angles 1 (at D) and 2 (at B).  6-29: triangle A B with apex B and
%   cevian BD to the base, angles 1 and 2 at D.
\newcommand{\FIGVITWENTYEIGHTTWENTYNINE}{\begin{bkfigure}[\linewidth]
\begin{minipage}{0.52\linewidth}\centering
\begin{tikzpicture}[scale=0.975]
  % Proportions off scan06 p21 (book p.91): the book's parallelogram is
  % width 1.8 x height, leaning right by 0.17 of its height.  The draft was
  % 2.3 x height and leaned 0.29, which flattened the plate enough to push
  % the two numerals into one another in the middle of the figure.
  \coordinate (A) at (0,0);
  \coordinate (B) at (2.62,0);
  \coordinate (D) at (0.25,1.45);
  \coordinate (C) at (2.87,1.45);
  \draw[given] (A) -- (B) -- (C) -- (D) -- cycle;
  \draw[given] (D) -- (B);
  \node[vlab, below left]  at (A) {$A$};
  \node[vlab, below right] at (B) {$B$};
  \node[vlab, above right] at (C) {$C$};
  \node[vlab, above left]  at (D) {$D$};
  \VIarc{C}{D}{B}{0.42cm}  \VInumb{C}{D}{B}{0.92cm}{1}
  \VIarc{A}{B}{D}{0.42cm}  \VInumb{A}{B}{D}{0.92cm}{2}
\end{tikzpicture}
\figcap{6--28}
\end{minipage}\hfill
\begin{minipage}{0.44\linewidth}\centering
\begin{tikzpicture}[scale=0.975]
  \coordinate (A) at (0,0);
  \coordinate (C) at (2.72,0);
  % Exercise 15 gives AD = DC and angle 1 = angle 2.  Angles 1 and 2 are the
  % two angles the downward stub of BD makes with AC, and they are
  % supplementary, so congruent means each is a right angle: B must stand
  % directly OVER the midpoint D, exactly as the book draws it.  The draft put
  % B at 0.34 of the base, which made the plate contradict its own hypothesis.
  \coordinate (B) at (1.36,1.62);
  % Exercise 15 gives AD = DC: D is the MIDPOINT of AC, not a point 0.42
  % along it, which is what the figure used to show.
  \coordinate (D) at ($(A)!0.5!(C)$);
  \draw[given] (A) -- (C) -- (B) -- cycle;
  \draw[given] (B) -- (D);
  \node[vlab, left]  at (A) {$A$};
  \node[vlab, right] at (C) {$C$};
  \node[vlab, above] at (B) {$B$};
  \node[vlab, above right] at (D) {$D$};
  % In the book BD is carried a little way BELOW the base, and the two
  % numbered angles are the ones under the line, spanned by a single arc.
  \draw[given] (D) -- ($(D)!-0.34!(B)$);
  % ADC is a straight angle, so which way \VIarc would sweep is a coin toss:
  % strike this one explicitly, below the base, as the book prints it.
  \draw[fig, line width=0.55pt] ([shift=(180:0.40cm)]D)
       arc[start angle=180, end angle=360, radius=0.40cm];
  \VInum{D}{225}{0.62cm}{1}
  \VInum{D}{315}{0.62cm}{2}
\end{tikzpicture}
\figcap{6--29}
\end{minipage}
\end{bkfigure}}

% =========================================================== Figs 6-30, 6-31
% Book p.91.  6-30: quadrilateral A B C D with E and F interior on the
%   diagonals; angles 1..6.  6-31: quadrilateral A B C D with centre O.
\newcommand{\FIGVITHIRTYTHIRTYONE}{\begin{bkfigure}[\linewidth]
\begin{minipage}{0.58\linewidth}\centering
% The book draws ONE diagonal here, A--C, and hangs two segments off it:
% B--F and E--D, with F and E on AC.  (An earlier draft drew the second
% diagonal B--D instead and left BF and ED out altogether, which made the
% six numbered angles label things that were not in the picture.)
% Exercise *16 gives AE = CF, so E and F sit symmetrically about the midpoint
% of AC -- placed here by ratio, not by eye.
\begin{tikzpicture}[scale=1.046]
  \coordinate (A) at (0,0);
  \coordinate (D) at (3.30,0);
  \coordinate (B) at (0.75,1.45);
  \coordinate (C) at (4.05,1.45);
  \draw[given] (A) -- (D) -- (C) -- (B) -- cycle;
  \draw[given] (A) -- (C);
  \coordinate (F) at ($(A)!0.33!(C)$);
  \coordinate (E) at ($(A)!0.67!(C)$);
  \draw[given] (B) -- (F);
  \draw[given] (E) -- (D);
  \node[vlab, below left]  at (A) {$A$};
  \node[vlab, below right] at (D) {$D$};
  \node[vlab, above left]  at (B) {$B$};
  \node[vlab, above right] at (C) {$C$};
  % E above-left of its point as the book sets it; F sits low on AC, close
  % to the edge AD, so its letter goes up into the wedge between AC and
  % the segment F--B.
  \node[vlab, above left] at (E) {$E$};
  \node[vlab] at ($(F)+(60:0.38)$) {$F$};
  % Radii are set per wedge: the two slivers at A and C (only about 20 degrees
  % wide, because BC is parallel to AD) need the numeral pushed well out
  % before it clears both sides, while the open wedges at F and E do not.
  \VIarc{B}{A}{C}{0.36cm}   \VInumb{B}{A}{C}{0.78cm}{3}
  \VIarc{C}{A}{D}{0.52cm}   \VInumb{C}{A}{D}{1.30cm}{6}
  \VIarc{A}{F}{B}{0.26cm}   \VInumb{A}{F}{B}{0.44cm}{1}
  \VIarc{C}{E}{D}{0.24cm}   \VInumb{C}{E}{D}{0.42cm}{2}
  \VIarc{B}{C}{A}{0.42cm}   \VInumb{B}{C}{A}{1.30cm}{5}
  \VIarc{A}{C}{D}{0.30cm}   \VInumb{A}{C}{D}{0.55cm}{4}
\end{tikzpicture}
\figcap{6--30}
\end{minipage}\hfill
\begin{minipage}{0.42\linewidth}\centering
\begin{tikzpicture}[scale=1.360]
  % Exercise *17 gives AB = AD and BO = OD, so B and D must sit symmetrically
  % about the line AC: B and D share an x, and AC passes through their midpoint.
  % Proportions off the plate: AO is about 2.5 times OC, and OB = OD is about
  % 0.37 of AO -- the book's kite is much longer in the nose than a first
  % draft suggests.
  \coordinate (A) at (0,0.68);
  \coordinate (B) at (1.85,1.36);
  \coordinate (C) at (2.60,0.68);
  \coordinate (D) at (1.85,0);
  \draw[given] (A) -- (B) -- (C) -- (D) -- cycle;
  \draw[given] (B) -- (D);
  \coordinate (O) at (intersection of A--C and B--D);
  \draw[given] (A) -- (C);
  \node[vlab, left]  at (A) {$A$};
  \node[vlab, above] at (B) {$B$};
  \node[vlab, right] at (C) {$C$};
  \node[vlab, below] at (D) {$D$};
  \node[vlab, left] at ($(O)+(0.02,-0.20)$) {$O$};
  \VItk{A}{B}{1}
  \VItk{A}{D}{1}
  \VItk{O}{B}{2}
  \VItk{O}{D}{2}
  % The pair of congruent angles at A, marked as the book marks them.  These
  % arcs used to be omitted because the old pic-based mark wrapped to a full
  % circle whenever the sides were named the other way round; the hand-struck
  % \VIarc has no such failure mode.
  \VIarc{B}{A}{C}{0.85cm}
  \VIarc{C}{A}{D}{0.85cm}
  % OC is only 0.75 long, so a numeral 0.68 out along the 45-degree bisector
  % lands OUTSIDE the kite, across BC.  The book keeps 1 and 2 tucked in
  % beside the double ticks, about 0.19 of AO from the crossing.
  \VIarc{B}{O}{C}{0.18cm}   \VInum{O}{35}{0.36cm}{1}
  \VIarc{C}{O}{D}{0.18cm}   \VInum{O}{-35}{0.36cm}{2}
\end{tikzpicture}
\figcap{6--31}
\end{minipage}
\end{bkfigure}}

% ================================================================= Fig 6-32
% Book p.92.  The isosceles triangle labelled two ways at once, which is the
% whole trick of the Theorem 6-3 proof.
\newcommand{\FIGVITHIRTYTWO}{\begin{bkfigure}[\linewidth]
\begin{tikzpicture}[scale=1.235]
  \coordinate (A) at (1.42,1.62);
  \coordinate (B) at (0,0);
  \coordinate (C) at (2.84,0);
  \draw[given] (A) -- (B) -- (C) -- cycle;
  \node[vlab, above] at (A) {$A = A\VIpr$};
  \node[vlab, below left]  at (B) {$B = C\VIpr$};
  \node[vlab, below right] at (C) {$C = B\VIpr$};
\end{tikzpicture}
\figcap{6--32}
\end{bkfigure}}

% ==================================================== Figs 6-33, 6-34, 6-35
% Book p.92, Exercise Group 6-6.
%  6-33: triangle P Q R.
%  6-34: the kite A B C D with both diagonals drawn.
%  6-35: triangle A D B with apex C and cevian CD, angles 1 and 2.
\newcommand{\FIGVITHIRTYTHREETHIRTYFOURTHIRTYFIVE}{\begin{bkfigure}[\linewidth]
\begin{minipage}{0.46\linewidth}\centering
\begin{tikzpicture}[scale=0.961]
  \coordinate (Q) at (0,0);
  \coordinate (R) at (2.62,0);
  \coordinate (P) at (1.22,1.12);
  \draw[given] (Q) -- (R) -- (P) -- cycle;
  \node[vlab, left]  at (Q) {$Q$};
  \node[vlab, right] at (R) {$R$};
  \node[vlab, above] at (P) {$P$};
\end{tikzpicture}
\figcap{6--33}

\medskip
\begin{tikzpicture}[scale=0.961]
  \coordinate (A) at (1.32,2.02);
  \coordinate (B) at (0,1.32);
  \coordinate (C) at (2.64,1.32);
  \coordinate (D) at (1.32,0);
  \draw[given] (A) -- (B) -- (D) -- (C) -- cycle;
  \draw[given] (B) -- (C);
  \draw[given] (A) -- (D);
  \node[vlab, above] at (A) {$A$};
  \node[vlab, left]  at (B) {$B$};
  \node[vlab, right] at (C) {$C$};
  \node[vlab, below] at (D) {$D$};
\end{tikzpicture}
\figcap{6--34}
\end{minipage}\hfill
\begin{minipage}{0.46\linewidth}\centering
\begin{tikzpicture}[scale=0.961]
  \coordinate (A) at (0,0);
  \coordinate (B) at (2.82,0);
  \coordinate (C) at (1.32,1.72);
  \coordinate (D) at ($(A)!0.47!(B)$);
  \draw[given] (A) -- (B) -- (C) -- cycle;
  \draw[given] (C) -- (D);
  \node[vlab, left]  at (A) {$A$};
  \node[vlab, right] at (B) {$B$};
  \node[vlab, above] at (C) {$C$};
  \node[vlab, below] at (D) {$D$};
  \VIarc{A}{C}{D}{0.46cm}  \VInumb{A}{C}{D}{0.76cm}{1}
  \VIarc{D}{C}{B}{0.46cm}  \VInumb{D}{C}{B}{0.72cm}{2}
\end{tikzpicture}
\figcap{6--35}
\end{minipage}
\end{bkfigure}}

% =========================================================== Figs 6-36, 6-37
% Book p.93.  6-36: the "dart" A B C E D -- AD is the long spine, B interior.
%             6-37: A B C D E F -- two triangles meeting at C/D, angles 1, 2.
\newcommand{\FIGVITHIRTYSIXTHIRTYSEVEN}{\begin{bkfigure}[\linewidth]
\begin{minipage}{0.50\linewidth}\centering
\begin{tikzpicture}[scale=0.972]
  \coordinate (A) at (0,0.86);
  \coordinate (C) at (1.02,1.62);
  \coordinate (E) at (1.02,0);
  \coordinate (D) at (2.72,0.86);
  \coordinate (B) at ($(A)!0.42!(D)$);
  \draw[given] (A) -- (D);
  \draw[given] (C) -- (B) -- (E);
  \draw[given] (C) -- (D) -- (E);
  \node[vlab, left]  at (A) {$A$};
  \node[vlab, above] at (C) {$C$};
  \node[vlab, below] at (E) {$E$};
  \node[vlab, right] at (D) {$D$};
  \node[vlab, below left] at (B) {$B$};
\end{tikzpicture}
\figcap{6--36}
\end{minipage}\hfill
\begin{minipage}{0.46\linewidth}\centering
\begin{tikzpicture}[scale=0.972]
  \coordinate (A) at (0,0.72);
  \coordinate (B) at (0.92,1.82);
  \coordinate (F) at (2.92,0.72);
  \coordinate (E) at (1.92,-0.22);
  % AC = DF is given, so C and D sit symmetrically on A--F.
  \coordinate (C) at ($(A)!0.35!(F)$);
  \coordinate (D) at ($(A)!0.65!(F)$);
  \draw[given] (A) -- (F);
  \draw[given] (A) -- (B) -- (D);
  \draw[given] (C) -- (E) -- (F);
  \node[vlab, left]  at (A) {$A$};
  \node[vlab, above] at (B) {$B$};
  \node[vlab, right] at (F) {$F$};
  \node[vlab, below] at (E) {$E$};
  \node[vlab, below left]  at (C) {$C$};
  \node[vlab, above right] at (D) {$D$};
  % Angle 2 is at D (between DB and DA), NOT at B: the S.A.S. pair the
  % exercise builds is (AD, DB, angle 2) against (FC, CE, angle 1), so both
  % marked angles sit on the line A--F.  The book strikes an arc on each; the
  % draft put 2 at vertex B and left both arcs off.
  \VIarc{A}{D}{B}{0.34cm}   \VInumb{A}{D}{B}{0.80cm}{2}
  \VIarc{F}{C}{E}{0.34cm}   \VInumb{F}{C}{E}{0.80cm}{1}
\end{tikzpicture}
\figcap{6--37}
\end{minipage}
\end{bkfigure}}

% =========================================================== Figs 6-38, 6-39
% Book p.93.  6-38: triangle A B C with E and D on the base and cevians from B.
%             6-39: triangle R S T with U inside, joined to S and T.
\newcommand{\FIGVITHIRTYEIGHTTHIRTYNINE}{\begin{bkfigure}[\linewidth]
\begin{minipage}{0.50\linewidth}\centering
\begin{tikzpicture}[scale=0.963]
  \coordinate (A) at (0,0);
  \coordinate (C) at (3.22,0);
  % Exercise 7 gives AD = EC and BE = BD.  The first puts E and D
  % symmetrically about the midpoint of AC; the second then puts B over that
  % midpoint.  The draft had E at 0.36, D at 0.56 and B at 0.41 of the base,
  % so neither hypothesis held in the drawing.
  \coordinate (B) at (1.61,1.62);
  \coordinate (E) at ($(A)!0.42!(C)$);
  \coordinate (D) at ($(A)!0.58!(C)$);
  \draw[given] (A) -- (C) -- (B) -- cycle;
  \draw[given] (B) -- (E);
  \draw[given] (B) -- (D);
  \node[vlab, left]  at (A) {$A$};
  \node[vlab, right] at (C) {$C$};
  \node[vlab, above] at (B) {$B$};
  \node[vlab, below] at (E) {$E$};
  \node[vlab, below] at (D) {$D$};
\end{tikzpicture}
\figcap{6--38}
\end{minipage}\hfill
\begin{minipage}{0.46\linewidth}\centering
\begin{tikzpicture}[scale=0.963]
  \coordinate (S) at (0,0);
  \coordinate (T) at (2.72,0);
  % RS = RT and US = UT are both given, so R and U must stand on the
  % perpendicular bisector of ST -- x = 1.36, not 1.22.
  \coordinate (R) at (1.36,1.72);
  \coordinate (U) at (1.36,0.92);
  \draw[given] (S) -- (T) -- (R) -- cycle;
  \draw[given] (S) -- (U) -- (T);
  \node[vlab, left]  at (S) {$S$};
  \node[vlab, right] at (T) {$T$};
  \node[vlab, above] at (R) {$R$};
  \node[vlab, below] at (U) {$U$};
\end{tikzpicture}
\figcap{6--39}
\end{minipage}
\end{bkfigure}}

% ==================================================== Figs 6-40, 6-41, 6-42
% Book p.93.
%  6-40: triangle A B with apex C; D and E on the sides, cevians AE and BD,
%        angles 1 and 2 at D and E.
%  6-41: triangle A D with apex E; B, F, C on the base, cevians from E.
%  6-42: triangle A B with apex C; G, D, O, E, F interior, angles 1..4.
\newcommand{\FIGVIFORTYFORTYONEFORTYTWO}{\begin{bkfigure}[\linewidth]
\begin{minipage}{0.46\linewidth}\centering
% Measured off the plate on book p.93 (photo PDF p.107): base 645 px, apex over
% the midpoint (AC = BC, as Exercise 9 requires), height/base = 0.674, D and E
% about 0.60 of the way up their sides.  The book sets D level with its point
% on the LEFT and E level with its point on the RIGHT -- not above them -- and
% leaves a clear digit's width between the numerals 1 and 2.
% Drawn at 1.05 rather than 0.80: D and E are only 1.17 units apart, and the
% two bisector-placed numerals aim at each other, so at the old scale the
% numerals had to sit so close to their vertices that both grazed the cevians
% (0.27 pt).  Widening the plate buys the clearance without moving either
% numeral off its own bisector.
\begin{tikzpicture}[scale=0.927]
  \coordinate (A) at (0,0);
  \coordinate (B) at (2.92,0);
  \coordinate (C) at (1.46,1.97);
  \coordinate (D) at ($(A)!0.60!(C)$);
  \coordinate (E) at ($(B)!0.60!(C)$);
  \draw[given] (A) -- (B) -- (C) -- cycle;
  \draw[given] (A) -- (E);
  \draw[given] (B) -- (D);
  \node[vlab, left]  at (A) {$A$};
  \node[vlab, right] at (B) {$B$};
  \node[vlab, above] at (C) {$C$};
  \node[vlab, left]  at (D) {$D$};
  \node[vlab, right] at (E) {$E$};
  % Angle 1 is CDB and angle 2 is CEA -- the pair the exercise asks about.
  % Both go on their own bisector, which is what keeps them off DB and EA.
  % The two bisectors point at EACH OTHER, so distance along them is what sets
  % the gap between the numerals: at 0.48/0.46 they closed to 1.2pt and read as
  % the single number "12".  Pulled back to 0.40 for the daylight the book
  % leaves (about one digit width); still clear of the 0.26 arcs.
  \VIarc{C}{D}{B}{0.26cm}   \VInumb{C}{D}{B}{0.40cm}{1}
  \VIarc{C}{E}{A}{0.26cm}   \VInumb{C}{E}{A}{0.40cm}{2}
\end{tikzpicture}
\figcap{6--40}

\medskip
% Rebuilt against scan07 p1 (book p.93) at 150 dpi, cross-read with the photo
% plate.  The earlier draft fanned every cevian from the APEX E, which cannot
% satisfy Exercise 10 at all: its hypotheses are AO = OD, AF = FD, BO = OC,
% BF = FC, so O and F are both distinct interior points joined to A, B, C, D.
% The plate agrees -- at A and at D two lines leave the vertex (one to E, one
% to O), and six lines meet at O (up to E, down to A, B, F, C, D).
% Structure: base A-B-F-C-D collinear; F the MIDPOINT of AD (AF = FD) and of
% BC (BF = FC); E directly above F, so the whole plate is mirror-symmetric
% about FE and AO = OD, BO = OC hold exactly rather than by eye.
% Proportions off the scan: base 825 px, F at 0.488 (drawn 0.5, as the
% exercise demands), B at 0.228 and C at 0.737 (drawn symmetric at
% 0.245/0.755), O 0.40 of the base above it, E 0.70 of the base above it.
\begin{tikzpicture}[scale=1.131]
  \coordinate (A) at (0,0);
  \coordinate (D) at (3.22,0);
  \coordinate (F) at ($(A)!0.5!(D)$);      % AF = FD
  \coordinate (B) at ($(A)!0.245!(D)$);
  \coordinate (C) at ($(A)!0.755!(D)$);    % BF = FC by symmetry about F
  \coordinate (E) at (1.61,2.254);         % directly above F
  % Re-measured on the plate: O sits 0.636 of the way from F up to E, not 0.593.
  \coordinate (O) at ($(F)!0.636!(E)$);    % on FE, so AO = OD and BO = OC
  \draw[given] (A) -- (D) -- (E) -- cycle;
  \draw[given] (E) -- (F);
  \draw[given] (O) -- (A);
  \draw[given] (O) -- (B);
  \draw[given] (O) -- (C);
  \draw[given] (O) -- (D);
  \node[vlab, left]  at (A) {$A$};
  \node[vlab, right] at (D) {$D$};
  \node[vlab, above] at (E) {$E$};
  \node[vlab, below] at (B) {$B$};
  \node[vlab, below] at (F) {$F$};
  \node[vlab, below] at (C) {$C$};
  % Six lines meet at O; the only clear quadrant is up-and-right, between OE
  % and OD, which is exactly where the book sets the letter.  That quadrant is
  % a wedge closing toward E -- bounded on the left by the vertical EF and on
  % the right by the side ED -- so the letter has to hug the vertical and stay
  % LOW, which is exactly what the plate does (its O starts about 2.5% of the
  % base to the right of the axis).  The default 4.6pt outer sep pushed the
  % box out along the 45-degree diagonal instead, straight into ED.
  \node[vlab, above right, inner sep=1pt, outer sep=2.4pt] at (O) {$O$};
\end{tikzpicture}
\figcap{6--41}
\end{minipage}\hfill
\begin{minipage}{0.50\linewidth}\centering
% Rebuilt against scan07 p1 (book p.93) at 200 dpi.  The cevians the book
% draws are A--F and B--G (F on CB, G on CA), crossing at O.  D and E are NOT
% on the base -- an earlier draft put them there, which is the one structural
% error the plate contradicts outright: the book marks D on the cevian A--F
% and E on the cevian B--G, and joins each back to C.  Measured off the plate,
% each sits 0.40 of the way from its base vertex along that cevian.
% C is placed exactly over the midpoint of AB so that AC = BC holds, and G, F
% are taken at the same ratio down the two sides so that CG = CF holds -- both
% are hypotheses of Exercise 11.
% Angles 1 and 2 sit at G and F; 3 and 4 where the cevians CD and CE cut
% BG and AF, which is where the book strikes their arcs.
% Drawn a little larger than the neighbouring plates: four numbered angles
% share the narrow band between G, P, Q and F, and at the smaller scale their
% bisector labels ran into one another in the middle of the triangle.
\begin{tikzpicture}[scale=1.148]
  \coordinate (A) at (0,0);
  \coordinate (B) at (4.30,0);
  \coordinate (C) at (2.15,2.64);
  \coordinate (G) at ($(C)!0.45!(A)$);
  \coordinate (F) at ($(C)!0.45!(B)$);
  \draw[given] (A) -- (B) -- (C) -- cycle;
  \draw[given] (A) -- (F);
  \draw[given] (B) -- (G);
  % Measured off scan07 p1: D and E sit 0.47 of the way along their cevians,
  % not 0.40.
  \coordinate (D) at ($(A)!0.47!(F)$);
  \coordinate (E) at ($(B)!0.47!(G)$);
  \draw[given] (C) -- (D);
  \draw[given] (C) -- (E);
  \coordinate (O) at (intersection of A--F and B--G);
  \coordinate (P) at (intersection of C--D and B--G);
  \coordinate (Q) at (intersection of C--E and A--F);
  \node[vlab, left]  at (A) {$A$};
  \node[vlab, right] at (B) {$B$};
  \node[vlab, above] at (C) {$C$};
  \node[vlab, below]  at (D) {$D$};
  \node[vlab, below] at (E) {$E$};
  \node[vlab, above left]  at (G) {$G$};
  \node[vlab, above right] at (F) {$F$};
  \node[vlab, below] at (O) {$O$};
  \VIarc{C}{G}{B}{0.20cm}   \VInumb{C}{G}{B}{0.38cm}{1}
  \VIarc{C}{F}{A}{0.20cm}   \VInumb{C}{F}{A}{0.38cm}{2}
  % Same trap as 6-40, and worse here because scale=1.30 multiplies the cm
  % distance too: the bisectors at P and Q aim straight at one another, so
  % 0.48cm drove both numerals onto the figure's axis and they printed as "34"
  % right on top of O.  0.36cm restores the gap the plate shows between them.
  \VIarc{C}{P}{B}{0.22cm}   \VInumb{C}{P}{B}{0.36cm}{3}
  \VIarc{C}{Q}{A}{0.22cm}   \VInumb{C}{Q}{A}{0.36cm}{4}
\end{tikzpicture}
\figcap{6--42}
\end{minipage}
\end{bkfigure}}

% ================================================================= Fig 6-43
% Book p.94.  The Theorem 6-4 proof picture: triangle ABC, and A'B'C' with the
% constructed point D' on ray C'B' (dashed, since it is added in the proof).
% The plate sets all four base-vertex letters BELOW the base line -- A and C
% under the first triangle, A' and C' under the second.  The draft lifted C and
% A' ABOVE it, which put them at the same height in the gap between the two
% triangles, where they read as one floating pair "C A'" belonging to neither.
% Panels also pulled 0.4cm further apart for the daylight the plate shows.
\newcommand{\FIGVIFORTYTHREE}{\begin{bkfigure}[\linewidth]
\begin{tikzpicture}[scale=1.235]
  \coordinate (A) at (0,0);
  \coordinate (C) at (2.52,0);
  \coordinate (B) at (1.02,1.42);
  \draw[given] (A) -- (C) -- (B) -- cycle;
  \VItk{A}{B}{1}  \VItk{A}{C}{2}  \VIarc{C}{A}{B}{0.46cm}
  \node[vlab, below left]  at (A) {$A$};
  \node[vlab, above] at (B) {$B$};
  \node[vlab, below right] at (C) {$C$};
  \begin{scope}[xshift=4.05cm]
   \coordinate (Ap) at (0,0);
   \coordinate (Cp) at (2.52,0);
   \coordinate (Bp) at (1.02,1.42);
   \coordinate (Dp) at ($(Cp)!0.30!(Bp)$);
   \draw[given] (Ap) -- (Cp) -- (Bp) -- cycle;
   \draw[aux]   (Ap) -- (Dp);
   \VItk{Ap}{Bp}{1}  \VItk{Ap}{Cp}{2}  \VIarc{Cp}{Ap}{Bp}{0.46cm}
   \node[vlab, below left]  at (Ap) {$A\VIpr$};
   \node[vlab, above] at (Bp) {$B\VIpr$};
   \node[vlab, below right] at (Cp) {$C\VIpr$};
   \node[vlab, above right] at (Dp) {$D\VIpr$};
  \end{scope}
\end{tikzpicture}
\figcap{6--43}
\end{bkfigure}}

% =========================================================== Figs 6-44, 6-45
% Book p.95, Exercise Group 6-8.
%  6-44: the rhombus A B C D with diagonal BD, angles 1 and 2 at B.
%  6-45: the parallelogram E F H G with diagonal E--H, angles 1 and 2.
\newcommand{\FIGVIFORTYFOURFORTYFIVE}{\begin{bkfigure}[\linewidth]
\begin{minipage}{0.46\linewidth}\centering
\begin{tikzpicture}[scale=0.982]
  % Off scan07 p3 the book's rhombus is very nearly square-on-its-corner
  % (half-width 117 against half-height 130); the draft was half again as
  % wide as it was tall.
  \coordinate (B) at (1.17,2.28);
  \coordinate (A) at (0,1.08);
  \coordinate (C) at (2.34,1.08);
  \coordinate (D) at (1.17,0);
  \draw[given] (A) -- (B) -- (C) -- (D) -- cycle;
  \draw[given] (B) -- (D);
  \node[vlab, above] at (B) {$B$};
  \node[vlab, left]  at (A) {$A$};
  \node[vlab, right] at (C) {$C$};
  \node[vlab, below] at (D) {$D$};
  \VInumb{A}{B}{D}{0.56cm}{1}
  \VInumb{D}{B}{C}{0.54cm}{2}
\end{tikzpicture}
\figcap{6--44}
\end{minipage}\hfill
\begin{minipage}{0.50\linewidth}\centering
\begin{tikzpicture}[scale=0.982]
  \coordinate (G) at (0,0);
  \coordinate (H) at (2.62,0);
  \coordinate (E) at (0.52,1.12);
  \coordinate (F) at (3.14,1.12);
  \draw[given] (G) -- (H) -- (F) -- (E) -- cycle;
  \draw[given] (E) -- (H);
  \node[vlab, below left]  at (G) {$G$};
  \node[vlab, below right] at (H) {$H$};
  \node[vlab, above left]  at (E) {$E$};
  \node[vlab, above right] at (F) {$F$};
  \VInumb{G}{E}{H}{0.38cm}{1}
  \VInumb{F}{H}{E}{0.38cm}{2}
\end{tikzpicture}
\figcap{6--45}
\end{minipage}
\end{bkfigure}}

% =========================================================== Figs 6-46, 6-47
% Book p.95.  6-46: triangle C D with apex A and cevian AB to the base.
%             6-47: triangle B D with apex A and cevian AC to the base.
\newcommand{\FIGVIFORTYSIXFORTYSEVEN}{\begin{bkfigure}[\linewidth]
\begin{minipage}{0.48\linewidth}\centering
\begin{tikzpicture}[scale=0.983]
  \coordinate (C) at (0,0);
  \coordinate (D) at (2.72,0);
  \coordinate (A) at (1.36,1.62);
  \coordinate (B) at ($(C)!0.5!(D)$);
  \draw[given] (C) -- (D) -- (A) -- cycle;
  \draw[given] (A) -- (B);
  \node[vlab, left]  at (C) {$C$};
  \node[vlab, right] at (D) {$D$};
  \node[vlab, above] at (A) {$A$};
  \node[vlab, below] at (B) {$B$};
\end{tikzpicture}
\figcap{6--46}
\end{minipage}\hfill
\begin{minipage}{0.48\linewidth}\centering
\begin{tikzpicture}[scale=0.983]
  \coordinate (B) at (0,0);
  \coordinate (D) at (2.72,0);
  % Exercise 4 gives AB = AD and BC = CD, so C is the MIDPOINT of BD and A
  % stands directly over it -- which is how the book draws it.  The draft had
  % C at 0.46 and A off to the left of it.
  \coordinate (C) at ($(B)!0.5!(D)$);
  \coordinate (A) at (1.36,1.62);
  \draw[given] (B) -- (D) -- (A) -- cycle;
  \draw[given] (A) -- (C);
  \node[vlab, left]  at (B) {$B$};
  \node[vlab, right] at (D) {$D$};
  \node[vlab, above] at (A) {$A$};
  \node[vlab, below] at (C) {$C$};
\end{tikzpicture}
\figcap{6--47}
\end{minipage}
\end{bkfigure}}

% =========================================================== Figs 6-48, 6-49
% Book p.96.  6-48: triangle A B C with D, E on the sides and both cevians
%             BE, CD drawn.  6-49: the kite B A D C with both diagonals.
\newcommand{\FIGVIFORTYEIGHTFORTYNINE}{\begin{bkfigure}[\linewidth]
\begin{minipage}{0.52\linewidth}\centering
\begin{tikzpicture}[scale=1.027]
  \coordinate (B) at (0,0);
  \coordinate (C) at (3.22,0);
  \coordinate (A) at (1.61,1.62);
  % Exercise 5 makes D and E the MIDPOINTS of AB and AC.
  \coordinate (D) at ($(B)!0.5!(A)$);
  \coordinate (E) at ($(C)!0.5!(A)$);
  \draw[given] (B) -- (C) -- (A) -- cycle;
  \draw[given] (B) -- (E);
  \draw[given] (C) -- (D);
  \draw[given] (D) -- (E);
  \node[vlab, left]  at (B) {$B$};
  \node[vlab, right] at (C) {$C$};
  \node[vlab, above] at (A) {$A$};
  \node[vlab, above left]  at (D) {$D$};
  \node[vlab, above right] at (E) {$E$};
\end{tikzpicture}
\figcap{6--48}
\end{minipage}\hfill
\begin{minipage}{0.42\linewidth}\centering
\begin{tikzpicture}[scale=1.027]
  % Off scan07 p3 the book's kite is TALL and narrow: half-width 0.85 of the
  % short (upper) axis, against a lower axis twice the upper.  The draft was
  % nearly twice as wide as the upper axis was long.
  \coordinate (A) at (1.22,1.92);
  \coordinate (B) at (0.68,1.28);
  \coordinate (D) at (1.76,1.28);
  \coordinate (C) at (1.22,0);
  \draw[given] (B) -- (A) -- (D) -- (C) -- cycle;
  \draw[given] (B) -- (D);
  \draw[given] (A) -- (C);
  \node[vlab, above] at (A) {$A$};
  \node[vlab, left]  at (B) {$B$};
  \node[vlab, right] at (D) {$D$};
  \node[vlab, below] at (C) {$C$};
\end{tikzpicture}
\figcap{6--49}
\end{minipage}
\end{bkfigure}}

% =========================================================== Figs 6-50, 6-51
% Book p.96.  6-50: triangle A B C, D and E midpoints of AB and BC, with the
%   crossing cevians AE and CD.  6-51: A B C on a line with peaks E and D.
\newcommand{\FIGVIFIFTYFIFTYONE}{\begin{bkfigure}[\linewidth]
\begin{minipage}{0.48\linewidth}\centering
\begin{tikzpicture}[scale=0.960]
  \coordinate (A) at (0,0);
  \coordinate (C) at (2.92,0);
  % Exercise 7 gives AB = BC, so B stands over the midpoint of AC; the draft
  % had it at 0.36 of the base and the drawing was not isosceles at all.
  \coordinate (B) at (1.46,1.62);
  \coordinate (D) at ($(A)!0.5!(B)$);
  \coordinate (E) at ($(B)!0.5!(C)$);
  \draw[given] (A) -- (C) -- (B) -- cycle;
  \draw[given] (A) -- (E);
  \draw[given] (C) -- (D);
  \node[vlab, left]  at (A) {$A$};
  \node[vlab, right] at (C) {$C$};
  \node[vlab, above] at (B) {$B$};
  \node[vlab, left]  at (D) {$D$};
  \node[vlab, above right] at (E) {$E$};
\end{tikzpicture}
\figcap{6--50}
\end{minipage}\hfill
\begin{minipage}{0.48\linewidth}\centering
\begin{tikzpicture}[scale=0.960]
  \coordinate (A) at (0,0);
  \coordinate (C) at (3.02,0);
  % Exercise 8 makes B the midpoint of AC and gives EB = BD, so E and D must
  % sit symmetrically about B.  The draft's E and D were 0.76 and 0.50 out.
  \coordinate (B) at ($(A)!0.5!(C)$);
  \coordinate (E) at (0.88,1.22);
  \coordinate (D) at (2.14,1.22);
  \draw[given] (A) -- (C);
  \draw[given] (A) -- (E) -- (B);
  \draw[given] (B) -- (D) -- (C);
  \node[vlab, left]  at (A) {$A$};
  \node[vlab, below] at (B) {$B$};
  \node[vlab, right] at (C) {$C$};
  \node[vlab, above] at (E) {$E$};
  \node[vlab, above] at (D) {$D$};
\end{tikzpicture}
\figcap{6--51}
\end{minipage}
\end{bkfigure}}

% ==================================================== Figs 6-52, 6-53, 6-54
% Book p.96.
%  6-52: quadrilateral D A B C with the crossing segments DB and AC.
%  6-53: triangle A B C with E, D on the base and cevians from B.
%  6-54: the farmer/barn word problem -- rectangle "Barn" with the dashed
%        sight-line A..B above and the crossed lines A--E, B--D through C.
\newcommand{\FIGVIFIFTYTWOFIFTYTHREEFIFTYFOUR}{\begin{bkfigure}[\linewidth]
\begin{minipage}{0.46\linewidth}\centering
\begin{tikzpicture}[scale=0.923]
  % Exercise 9 gives angle DAB = angle CBA and AD = BC, so the quadrilateral
  % is an ISOSCELES trapezoid; the draft's D and C were 0.42 and 0.50 in from
  % the ends, which broke both hypotheses.
  \coordinate (A) at (0,0);
  \coordinate (B) at (3.02,0);
  \coordinate (D) at (0.46,1.22);
  \coordinate (C) at (2.56,1.22);
  \draw[given] (A) -- (B);
  \draw[given] (D) -- (C);
  \draw[given] (D) -- (A);
  \draw[given] (C) -- (B);
  \draw[given] (D) -- (B);
  \draw[given] (A) -- (C);
  \node[vlab, below left]  at (A) {$A$};
  \node[vlab, below right] at (B) {$B$};
  \node[vlab, above left]  at (D) {$D$};
  \node[vlab, above right] at (C) {$C$};
\end{tikzpicture}
\figcap{6--52}

\medskip
\begin{tikzpicture}[scale=0.923]
  % Exercise 10 gives AD = EC and BE = BD, the same pair of hypotheses as
  % Fig. 6-38: E and D symmetric about the midpoint, B over that midpoint.
  \coordinate (A) at (0,0);
  \coordinate (C) at (3.02,0);
  \coordinate (B) at (1.51,1.62);
  \coordinate (E) at ($(A)!0.42!(C)$);
  \coordinate (D) at ($(A)!0.58!(C)$);
  \draw[given] (A) -- (C) -- (B) -- cycle;
  \draw[given] (B) -- (E);
  \draw[given] (B) -- (D);
  \node[vlab, left]  at (A) {$A$};
  \node[vlab, right] at (C) {$C$};
  \node[vlab, above] at (B) {$B$};
  \node[vlab, below] at (E) {$E$};
  \node[vlab, below] at (D) {$D$};
\end{tikzpicture}
\figcap{6--53}
\end{minipage}\hfill
\begin{minipage}{0.50\linewidth}\centering
\begin{tikzpicture}[scale=0.923]
  \coordinate (A) at (0,1.72);
  \coordinate (B) at (3.42,1.72);
  \coordinate (D) at (0.42,0);
  \coordinate (E) at (3.02,0);
  \draw[given] (A) -- (E);
  \draw[given] (B) -- (D);
  \coordinate (C) at (intersection of A--E and B--D);
  % the book's barn is drawn round the word with clear space; at the old
  % height the label was as tall as the box and touched both long sides.
  \draw[given] (1.02,1.42) rectangle (2.42,2.02);
  \node[slab] at (1.72,1.72) {Barn};
  \draw[aux] (A) -- (1.02,1.72);
  \draw[aux] (2.42,1.72) -- (B);
  \node[vlab, left]  at (A) {$A$};
  \node[vlab, right] at (B) {$B$};
  \node[vlab, left]  at (D) {$D$};
  \node[vlab, right] at (E) {$E$};
  \node[vlab] at ($(C)+(0:0.66)$) {$C$};
\end{tikzpicture}
\figcap{6--54}
\end{minipage}
\end{bkfigure}}

% ================================================================= Fig 6-55
% Book p.97.  The A.S.A. proof picture -- triangles ABC and A'B'C' with
% the constructed D' on ray A'B' (dashed), arcs at A/A' and C/C', tick on
% AC/A'C'.
% Fig 6-56 used to be bundled into this block and emitted beside Theorem 6-5.
% It belongs to Exercise Group 6-9 #1, further down p.97, and is now its own
% macro so it can be placed there.
\newcommand{\FIGVIFIFTYFIVE}{\begin{bkfigure}[\linewidth]
\begin{tikzpicture}[scale=1.196]
  \coordinate (A) at (0,0);
  \coordinate (C) at (2.62,0);
  \coordinate (B) at (0.62,1.42);
  \draw[given] (A) -- (C) -- (B) -- cycle;
  \VItk{A}{C}{1}
  \VIarc{C}{A}{B}{0.46cm}
  \VIarcb{B}{C}{A}{0.46cm}
  \node[vlab, below left]  at (A) {$A$};
  \node[vlab, above] at (B) {$B$};
  % C and A' face each other across the gap between the two triangles.  Set
  % above the base line and with only 1.34cm of canvas between the panels they
  % closed up and printed as one word, "CA'".  Both base vertices now take the
  % same below-the-line placement A and C' already use, and the panels are
  % 4.10cm apart, which is the daylight the book leaves.
  \node[vlab, below right] at (C) {$C$};
  \begin{scope}[xshift=4.10cm]
   \coordinate (Ap) at (0,0);
   \coordinate (Cp) at (2.62,0);
   \coordinate (Bp) at (0.62,1.42);
   \coordinate (Dp) at ($(Ap)!0.72!(Bp)$);
   \draw[given] (Ap) -- (Cp) -- (Bp) -- cycle;
   \draw[aux]   (Dp) -- (Cp);
   \VItk{Ap}{Cp}{1}
   \VIarc{Cp}{Ap}{Bp}{0.46cm}
   \VIarcb{Bp}{Cp}{Ap}{0.46cm}
   \node[vlab, below left]  at (Ap) {$A\VIpr$};
   \node[vlab, above] at (Bp) {$B\VIpr$};
   \node[vlab, below right] at (Cp) {$C\VIpr$};
   \node[vlab, left] at (Dp) {$D\VIpr$};
  \end{scope}
\end{tikzpicture}
\figcap{6--55}
\end{bkfigure}}

% ================================================================= Fig 6-56
% Book p.97, Exercise Group 6-9 #1.  The kite A B C D with the spine B--D:
% angles 1 and 2 under B, angles 3 and 4 over D.  Checked against scan07 p6
% (book p.97): the book strikes an ARC across each of the four angles, which
% the earlier draft left off, and every numeral sits inside its own arc on
% that angle's bisector.
\newcommand{\FIGVIFIFTYSIX}{\begin{bkfigure}[\linewidth]
\begin{tikzpicture}[scale=1.118]
  \coordinate (B) at (1.32,1.62);
  \coordinate (D) at (1.32,0);
  \coordinate (A) at (0,1.02);
  \coordinate (C) at (2.64,1.02);
  \draw[given] (A) -- (B) -- (C);
  \draw[given] (A) -- (D) -- (C);
  \draw[given] (B) -- (D);
  \node[vlab, above] at (B) {$B$};
  \node[vlab, below] at (D) {$D$};
  \node[vlab, left]  at (A) {$A$};
  \node[vlab, right] at (C) {$C$};
  \VIarc{A}{B}{D}{0.34cm}   \VInumb{A}{B}{D}{0.54cm}{1}
  \VIarc{D}{B}{C}{0.34cm}   \VInumb{D}{B}{C}{0.54cm}{2}
  \VIarc{A}{D}{B}{0.34cm}   \VInumb{A}{D}{B}{0.62cm}{3}
  \VIarc{B}{D}{C}{0.34cm}   \VInumb{B}{D}{C}{0.62cm}{4}
\end{tikzpicture}
\figcap{6--56}
\end{bkfigure}}

% =========================================================== Figs 6-57, 6-58
% Book p.98.  6-57: parallelogram A D C B with diagonal B--D, angles 1..4.
%             6-58: triangle A C with apex B, E and D on the base, angles 1, 2.
\newcommand{\FIGVIFIFTYSEVENFIFTYEIGHT}{\begin{bkfigure}[\linewidth]
\begin{minipage}{0.50\linewidth}\centering
% Rebuilt against scan07 p8 (book p.98).  The book draws ONE diagonal, B--D,
% and all four numbered angles are the parts it cuts off: 3 = ABD and
% 1 = DBC at B, 2 = ADB and 4 = BDC at D, each carrying its own arc.  The
% earlier draft drew a second diagonal A--C that is not in the plate, put 3 at
% vertex A, and stacked 2 and 4 on the SAME bisector at two radii -- which is
% why the collision audit read them as one word "42".
\begin{tikzpicture}[scale=1.289]
  \coordinate (B) at (0,0);
  \coordinate (C) at (2.92,0);
  \coordinate (A) at (0.42,1.22);
  \coordinate (D) at (3.34,1.22);
  \draw[given] (B) -- (C) -- (D) -- (A) -- cycle;
  \draw[given] (B) -- (D);
  \node[vlab, below left]  at (B) {$B$};
  \node[vlab, below right] at (C) {$C$};
  \node[vlab, above left]  at (A) {$A$};
  \node[vlab, above right] at (D) {$D$};
  \VIarc{A}{B}{D}{0.52cm}   \VInumb{A}{B}{D}{0.86cm}{3}
  \VIarc{D}{B}{C}{0.34cm}   \VInumb{D}{B}{C}{1.05cm}{1}
  \VIarc{A}{D}{B}{0.34cm}   \VInumb{A}{D}{B}{1.05cm}{2}
  \VIarc{B}{D}{C}{0.52cm}   \VInumb{B}{D}{C}{0.86cm}{4}
\end{tikzpicture}
\figcap{6--57}
\end{minipage}\hfill
\begin{minipage}{0.46\linewidth}\centering
\begin{tikzpicture}[scale=0.881]
  \coordinate (A) at (0,0);
  \coordinate (C) at (3.02,0);
  \coordinate (B) at (1.32,1.72);
  \coordinate (E) at ($(A)!0.30!(C)$);
  \coordinate (D) at ($(A)!0.62!(C)$);
  \draw[given] (A) -- (C) -- (B) -- cycle;
  \draw[given] (B) -- (E);
  \draw[given] (B) -- (D);
  \node[vlab, left]  at (A) {$A$};
  \node[vlab, right] at (C) {$C$};
  \node[vlab, above] at (B) {$B$};
  \node[vlab, below] at (E) {$E$};
  \node[vlab, below] at (D) {$D$};
  \VInumb{C}{B}{A}{0.88cm}{1}
  \VInumb{D}{B}{C}{1.10cm}{2}
\end{tikzpicture}
\figcap{6--58}
\end{minipage}
\end{bkfigure}}

% ==================================================== Figs 6-59, 6-60
% Book p.98.  6-59: the "dart" A B C E D with the spine through A and B.
%             6-60: triangle A B C with E, D on the sides and crossing cevians.
\newcommand{\FIGVIFIFTYNINESIXTY}{\begin{bkfigure}[\linewidth]
\begin{minipage}{0.48\linewidth}\centering
\begin{tikzpicture}[scale=1.273]
  \coordinate (A) at (0.32,0.92);
  \coordinate (B) at (1.32,0.92);
  \coordinate (C) at (0.92,1.72);
  \coordinate (E) at (0.92,0.10);
  \coordinate (D) at (2.62,0.92);
  \draw[given] (0,0.92) -- (D);
  \draw[given] (C) -- (B) -- (E);
  \draw[given] (C) -- (D) -- (E);
  \node[vlab, above left]  at (A) {$A$};
  % B's letter used to sit up-and-left, straight along C--B; up-and-right
  % runs into C--D, so it drops into the band between the spine and E--D.
  \node[vlab] at ($(B)+(300:0.32)$) {$B$};
  \node[vlab, above] at (C) {$C$};
  \node[vlab, below] at (E) {$E$};
  \node[vlab, right] at (D) {$D$};
\end{tikzpicture}
\figcap{6--59}
\end{minipage}\hfill
\begin{minipage}{0.48\linewidth}\centering
\begin{tikzpicture}[scale=0.926]
  \coordinate (A) at (0,0);
  \coordinate (C) at (3.12,0);
  \coordinate (B) at (1.42,1.62);
  \coordinate (E) at ($(A)!0.52!(B)$);
  \coordinate (D) at ($(C)!0.52!(B)$);
  \draw[given] (A) -- (C) -- (B) -- cycle;
  \draw[given] (A) -- (D);
  \draw[given] (C) -- (E);
  \node[vlab, left]  at (A) {$A$};
  \node[vlab, right] at (C) {$C$};
  \node[vlab, above] at (B) {$B$};
  \node[vlab, above left]  at (E) {$E$};
  \node[vlab, above right] at (D) {$D$};
\end{tikzpicture}
\figcap{6--60}
\end{minipage}
\end{bkfigure}}

% ==================================================== Figs 6-61, 6-62, 6-63
% Book p.98.
%  6-61: segments A--D and E--C crossing at B.
%  6-62: the "kite" A F C D E -- F and C on the spine A--D, E below.
%  6-63: parallelogram-ish B C A D with the base extended past D, angles 3, 4.
\newcommand{\FIGVISIXTYONESIXTYTWOSIXTYTHREE}{\begin{bkfigure}[\linewidth]
\begin{minipage}{0.46\linewidth}\centering
\begin{tikzpicture}[scale=0.914]
  \coordinate (A) at (0,1.22);
  \coordinate (D) at (3.02,0);
  \coordinate (E) at (0,0);
  \coordinate (C) at (3.02,1.22);
  \draw[given] (A) -- (D);
  \draw[given] (E) -- (C);
  \coordinate (B) at (intersection of A--D and E--C);
  \node[vlab, left]  at (A) {$A$};
  \node[vlab, right] at (D) {$D$};
  \node[vlab, left]  at (E) {$E$};
  \node[vlab, right] at (C) {$C$};
  \node[vlab, above] at (B) {$B$};
\end{tikzpicture}
\figcap{6--61}

\medskip
\begin{tikzpicture}[scale=1.256]
  \coordinate (A) at (0,0.72);
  \coordinate (D) at (3.22,0.72);
  \coordinate (B) at (1.12,1.62);
  \coordinate (E) at (1.72,0);
  \coordinate (F) at ($(A)!0.34!(D)$);
  \coordinate (C) at ($(A)!0.66!(D)$);
  \draw[given] (A) -- (D);
  \draw[given] (A) -- (B) -- (C);
  \draw[given] (F) -- (E) -- (D);
  \draw[given] (B) -- (D);
  \node[vlab, left]  at (A) {$A$};
  \node[vlab, right] at (D) {$D$};
  \node[vlab, above] at (B) {$B$};
  \node[vlab, below] at (E) {$E$};
  \node[vlab, below left]  at (F) {$F$};
  % C is the foot of B--C on the spine; below-LEFT lay the letter along
  % that segment, so it goes straight down into the band above E--D.
  \node[vlab, below] at (C) {$C$};
\end{tikzpicture}
\figcap{6--62}
\end{minipage}\hfill
\begin{minipage}{0.48\linewidth}\centering
\begin{tikzpicture}[scale=0.914]
  % Exercise 10 gives AB = AD and "AC bisects angle BAD", so B and D are
  % equidistant from A and C stands on the bisector of that angle.  The
  % draft's AB was 1.33 against AD 2.12 and its AC missed the bisector
  % altogether -- the plate showed neither hypothesis.
  \coordinate (A) at (0,0);
  \coordinate (D) at ($(A)+(0:1.55)$);
  \coordinate (B) at ($(A)+(70:1.55)$);
  \coordinate (C) at ($(A)+(35:2.30)$);
  \coordinate (X) at (2.30,0);          % the base carried on past D
  \draw[given] (A) -- (X);
  \draw[given] (A) -- (B) -- (C);
  \draw[given] (A) -- (C);
  \draw[given] (C) -- (D);
  \node[vlab, left]  at (A) {$A$};
  \node[vlab, below] at (D) {$D$};
  \node[vlab, above left]  at (B) {$B$};
  \node[vlab, above right] at (C) {$C$};
  % 4 is the angle between DC and the STUB of the base beyond D -- the book
  % extends the base for exactly that purpose.  The draft aimed the numeral
  % along DC itself.
  \VIarc{A}{B}{C}{0.34cm}   \VInumb{A}{B}{C}{0.58cm}{3}
  \VIarc{C}{D}{X}{0.30cm}   \VInumb{C}{D}{X}{0.54cm}{4}
\end{tikzpicture}
\figcap{6--63}
\end{minipage}
\end{bkfigure}}

% =========================================================== Figs 6-64, 6-65
% Book p.99.  6-64: triangle A B with apex C, D and E on the sides, crossing
%   cevians AE and BD meeting at F, angles 3 and 4 at D and E.
%  6-65: triangle A C with apex B and D below B, angles 1 and 2.
\newcommand{\FIGVISIXTYFOURSIXTYFIVE}{\begin{bkfigure}[\linewidth]
\begin{minipage}{0.50\linewidth}\centering
\begin{tikzpicture}[scale=0.935]
  \coordinate (A) at (0,0);
  \coordinate (B) at (3.22,0);
  \coordinate (C) at (1.61,1.72);
  \coordinate (D) at ($(A)!0.62!(C)$);
  \coordinate (E) at ($(B)!0.62!(C)$);
  \draw[given] (A) -- (B) -- (C) -- cycle;
  \draw[given] (A) -- (E);
  \draw[given] (B) -- (D);
  \coordinate (F) at (intersection of A--E and B--D);
  \node[vlab, left]  at (A) {$A$};
  \node[vlab, right] at (B) {$B$};
  \node[vlab, above] at (C) {$C$};
  \node[vlab, above left]  at (D) {$D$};
  \node[vlab, above right] at (E) {$E$};
  \node[vlab, below] at (F) {$F$};
  % 3 is angle CDB at D and 4 is angle CEA at E, each on its own bisector and
  % each carrying an arc.  The draft aimed them at fixed bearings that fell
  % straight along DB and EA.
  % The two bisectors aim at one another, so 0.58 drove the numerals together
  % until they printed as "34" in a single heap; 0.44 gives them the book's gap
  % and still clears the 0.26 arcs.
  \VIarc{C}{D}{B}{0.26cm}   \VInumb{C}{D}{B}{0.44cm}{3}
  \VIarc{C}{E}{A}{0.26cm}   \VInumb{C}{E}{A}{0.44cm}{4}
\end{tikzpicture}
\figcap{6--64}
\end{minipage}\hfill
\begin{minipage}{0.46\linewidth}\centering
\begin{tikzpicture}[scale=0.935]
  % Exercise 12 gives AB = BC and AD = DC, so B and D BOTH lie on the
  % perpendicular bisector of AC, which is how the book draws them -- D
  % directly under B, half its height.  The draft put B at 0.35 of the base
  % and D off the axis, so neither hypothesis held.
  \coordinate (A) at (0,0);
  \coordinate (C) at (2.92,0);
  \coordinate (B) at (1.46,1.82);
  \coordinate (D) at (1.46,0.91);
  \draw[given] (A) -- (C) -- (B) -- cycle;
  \draw[given] (A) -- (D) -- (C);
  \draw[given] (B) -- (D);
  \node[vlab, left]  at (A) {$A$};
  \node[vlab, right] at (C) {$C$};
  \node[vlab, above] at (B) {$B$};
  \node[vlab, below] at (D) {$D$};
  \VIarc{A}{B}{D}{0.40cm}   \VInumb{A}{B}{D}{0.72cm}{1}
  \VIarc{D}{B}{C}{0.40cm}   \VInumb{D}{B}{C}{0.72cm}{2}
\end{tikzpicture}
\figcap{6--65}
\end{minipage}
\end{bkfigure}}

% =========================================================== Figs 6-66, 6-67
% Book p.99.  Both are the "cevian web" figures.
%  6-66: triangle A B with apex C; D, E on the sides, H on AB; cevians AE, BD
%        and CH concurrent at G; angles 3, 4 at D, E.
%  6-67: triangle A B with apex C; D, E on the sides; F, G on AB; cevians
%        crossing at H, I, J; angles 1, 2 at the base.
\newcommand{\FIGVISIXTYSIXSIXTYSEVEN}{\begin{bkfigure}[\linewidth]
\begin{minipage}{0.48\linewidth}\centering
% Re-measured on the plate (book p.99, photo PDF p.113) at 4x: base 700 px,
% apex 610 px above it -- height/base 0.871, NOT the 0.565 the draft drew.
% Flattening the triangle is what squeezed the band between D--E and the
% cevian crossing G, and that is the band the letters F and G live in: F was
% pushed 13.7% of the base out from its point to find room and landed on A--E.
% D and E measured at 0.60 of the way up their sides (draft had 0.66).
\begin{tikzpicture}[scale=1.50]
  \coordinate (A) at (0,0);
  \coordinate (B) at (3.22,0);
  \coordinate (C) at (1.61,2.80);
  \coordinate (H) at ($(A)!0.5!(B)$);
  \coordinate (D) at ($(A)!0.60!(C)$);
  \coordinate (E) at ($(B)!0.60!(C)$);
  \draw[given] (A) -- (B) -- (C) -- cycle;
  \draw[given] (A) -- (E);
  \draw[given] (B) -- (D);
  \draw[given] (C) -- (H);
  \coordinate (G) at (intersection of A--E and B--D);
  \draw[given] (D) -- (E);
  % F is where C--H crosses D--E; the book letters it and the draft left it
  % out altogether.
  \coordinate (F) at (intersection of C--H and D--E);
  \node[vlab, left]  at (A) {$A$};
  \node[vlab, right] at (B) {$B$};
  \node[vlab, above] at (C) {$C$};
  \node[vlab, below] at (H) {$H$};
  \node[vlab, above left]  at (D) {$D$};
  \node[vlab, above right] at (E) {$E$};
  % Offsets measured off the plate, as a fraction of the base: F sits about
  % 5% out from its point and G about 9%, both down-and-left.  The draft used
  % 13.7% and 14.3%, which is what carried F across the cevian A--E.
  \node[vlab] at ($(F)+(225:0.227)$) {$F$};
  \node[vlab] at ($(G)+(247:0.33)$) {$G$};
  % 3 and 4 are the two halves of the vertex angle at C, and the book sets
  % them at the SAME radius on either side of C--H.  The draft used 1.30 and
  % 0.34, which threw both numerals onto the sides of the triangle.
  \VIarc{A}{C}{H}{0.30cm}   \VInumb{A}{C}{H}{0.44cm}{3}
  \VIarc{H}{C}{B}{0.30cm}   \VInumb{H}{C}{B}{0.44cm}{4}
\end{tikzpicture}
\figcap{6--66}
\end{minipage}\hfill
\begin{minipage}{0.48\linewidth}\centering
% Rebuilt against the photo of book p.99 (PDF 113).  F and G are on the SIDES
% CA and CB (below D and E), not on the base: the segment F--G runs across the
% triangle, and the two lines drawn are the cevians A--E and B--D.  H is where
% those cevians cross; I and J are where they cut F--G.  The earlier draft put
% F and G on AB, drew D--G and E--F instead, and then placed I and J at
% arbitrary ratios along F--G rather than deriving them, so the exercise's
% triangle HIJ was not the one in the picture.
% Exercise 14 gives AC = BC and DF = EG, and the plate is mirror-symmetric
% about the altitude from C -- which is what makes triangle HIJ isosceles.
\begin{tikzpicture}[scale=1.30]
  \coordinate (A) at (0,0);
  \coordinate (B) at (3.32,0);
  \coordinate (C) at (1.66,2.70);
  \coordinate (D) at ($(A)!0.55!(C)$);
  \coordinate (E) at ($(B)!0.55!(C)$);
  \coordinate (F) at ($(A)!0.30!(C)$);
  \coordinate (G) at ($(B)!0.30!(C)$);
  \draw[given] (A) -- (B) -- (C) -- cycle;
  \draw[given] (F) -- (G);
  \draw[given] (A) -- (E);
  \draw[given] (B) -- (D);
  \coordinate (H) at (intersection of A--E and B--D);
  \coordinate (I) at (intersection of A--E and F--G);
  \coordinate (J) at (intersection of B--D and F--G);
  \node[vlab, left]  at (A) {$A$};
  \node[vlab, right] at (B) {$B$};
  \node[vlab, above] at (C) {$C$};
  \node[vlab, above left]  at (D) {$D$};
  \node[vlab, above right] at (E) {$E$};
  \node[vlab, left]  at (F) {$F$};
  \node[vlab, right] at (G) {$G$};
  \node[vlab, above] at (H) {$H$};
  \node[vlab] at ($(I)+(285:0.32)$) {$I$};
  \node[vlab] at ($(J)+(255:0.32)$) {$J$};
  % 1 is the angle at A between the side AC and the cevian AE; 2 its mirror
  % image at B.
  \VIarc{C}{A}{E}{0.40cm}   \VInumb{C}{A}{E}{0.72cm}{1}
  \VIarc{C}{B}{D}{0.40cm}   \VInumb{C}{B}{D}{0.72cm}{2}
\end{tikzpicture}
\figcap{6--67}
\end{minipage}
\end{bkfigure}}

% ================================================================= Fig 6-68
% Book p.99.  Theorem 6-6: the same triangle labelled two ways, as in 6-32.
\newcommand{\FIGVISIXTYEIGHT}{\begin{bkfigure}[\linewidth]
\begin{tikzpicture}[scale=1.196]
  \coordinate (B) at (1.32,1.52);
  \coordinate (A) at (0,0);
  \coordinate (C) at (2.64,0);
  \draw[given] (A) -- (B) -- (C) -- cycle;
  \node[vlab, above] at (B) {$B = B\VIpr$};
  \node[vlab, below left]  at (A) {$A = C\VIpr$};
  \node[vlab, below right] at (C) {$C = A\VIpr$};
\end{tikzpicture}
\figcap{6--68}
\end{bkfigure}}

% ================================================================= Fig 6-69
% Book p.100.  The S.S.S. proof: (a) triangle ABC, (b) triangle A'B'C' with the
% arc at A', (c) the reflected copy with E and D below AC (ticks on AB, AE).
\newcommand{\FIGVISIXTYNINE}{\begin{bkfigure}[\linewidth]
\begin{tikzpicture}[scale=1.118]
  % (a)
  \coordinate (A) at (0,0);
  \coordinate (C) at (1.92,0);
  \coordinate (B) at (0.72,1.12);
  \draw[given] (A) -- (C) -- (B) -- cycle;
  \node[vlab, left]  at (A) {$A$};
  \node[vlab, below left] at (C) {$C$};
  \node[vlab, above] at (B) {$B$};
  \node[slab, below] at (0.96,-0.28) {(a)};
  % (b)
  \begin{scope}[xshift=2.72cm]
   \coordinate (Ap) at (0,0);
   \coordinate (Cp) at (1.92,0);
   \coordinate (Bp) at (0.72,1.12);
   \draw[given] (Ap) -- (Cp) -- (Bp) -- cycle;
   \VIarc{Cp}{Ap}{Bp}{0.42cm}
   \node[vlab, below]  at (Ap) {$A\VIpr$};
   % Set flush right, C' stacked under panel (c)'s A with 0.35pt between them.
   % Below-right matches the placement A' already takes at the other end of the
   % same base and drops it clear of the next panel.
   \node[vlab, below right] at (Cp) {$C\VIpr$};
   \node[vlab, above] at (Bp) {$B\VIpr$};
   \node[slab, below] at (0.96,-0.28) {(b)};
  \end{scope}
  % (c)
  \begin{scope}[xshift=5.44cm]
   \coordinate (Ac) at (0,0.52);
   \coordinate (Cc) at (1.72,0.52);
   \coordinate (Bc) at (1.22,1.52);
   \coordinate (Ec) at (1.22,-0.48);
   \coordinate (Dc) at (1.46,-0.62);
   \draw[given] (Ac) -- (Cc);
   \draw[given] (Ac) -- (Bc) -- (Cc);
   \draw[given] (Ac) -- (Ec) -- (Cc);
   \draw[given] (Ac) -- (Dc);
   \VItk{Ac}{Bc}{1}
   \VItk{Ac}{Ec}{1}
   \VIarc{Cc}{Ac}{Ec}{0.40cm}
   \node[vlab, left]  at (Ac) {$A$};
   \node[vlab, right] at (Cc) {$C$};
   \node[vlab, above] at (Bc) {$B$};
   \node[vlab, below]  at (Ec) {$E$};
   \node[vlab, right] at (Dc) {$D$};
   \node[slab, below] at (0.62,-0.92) {(c)};
  \end{scope}
\end{tikzpicture}
\figcap{6--69}
\end{bkfigure}}

% ================================================================= Fig 6-70
% Book p.100.  The kite A B C E with both diagonals, angles 1..4 at B and E.
\newcommand{\FIGVISEVENTY}{\begin{bkfigure}[\linewidth]
\begin{tikzpicture}[scale=1.196]
  \coordinate (B) at (1.32,1.72);
  \coordinate (A) at (0,0.86);
  \coordinate (C) at (2.64,0.86);
  \coordinate (E) at (1.32,-0.32);
  \draw[given] (A) -- (B) -- (C) -- (E) -- cycle;
  \draw[given] (A) -- (C);
  \draw[given] (B) -- (E);
  \node[vlab, above] at (B) {$B$};
  \node[vlab, left]  at (A) {$A$};
  \node[vlab, right] at (C) {$C$};
  \node[vlab, below] at (E) {$E$};
  \VInum{B}{-118}{0.42cm}{1}
  \VInum{B}{-62}{0.42cm}{2}
  \VInumb{B}{E}{A}{0.64cm}{3}
  \VInumb{C}{E}{B}{0.64cm}{4}
\end{tikzpicture}
\figcap{6--70}
\end{bkfigure}}

% ================================================================= Fig 6-71
% Book p.101.  The two cases the S.S.S. proof tacitly assumed away: BE meets AC
% inside the segment (a) or outside it (b).
\newcommand{\FIGVISEVENTYONE}{\begin{bkfigure}[\linewidth]
\begin{tikzpicture}[scale=1.144]
  % (a) C between A and the crossing
  \coordinate (A) at (0,0.62);
  \coordinate (B) at (1.42,1.52);
  \coordinate (E) at (1.42,-0.42);
  \coordinate (C) at (0.92,0.62);
  \draw[given] (A) -- (B) -- (E) -- cycle;
  \draw[given] (A) -- (C);
  \draw[given] (B) -- (E);
  \node[vlab, left]  at (A) {$A$};
  \node[vlab, above] at (B) {$B$};
  \node[vlab, below] at (E) {$E$};
  \node[vlab, right] at (C) {$C$};
  \node[slab, below] at (0.86,-0.86) {(a)};
  \begin{scope}[xshift=2.72cm]
   \coordinate (A2) at (0,0.62);
   \coordinate (B2) at (1.42,1.52);
   \coordinate (E2) at (1.42,-0.42);
   % C2 sat close enough to the vertical B2--E2 that its letter, set to the
   % right, ran into that segment.
   \coordinate (C2) at (0.88,0.62);
   \draw[given] (A2) -- (B2) -- (E2) -- cycle;
   \draw[given] (A2) -- (C2);
   \draw[given] (B2) -- (C2) -- (E2);
   \node[vlab, left]  at (A2) {$A$};
   \node[vlab, above] at (B2) {$B$};
   \node[vlab, below] at (E2) {$E$};
   \node[vlab, right] at (C2) {$C$};
   \node[slab, below] at (0.86,-0.86) {(b)};
  \end{scope}
\end{tikzpicture}
\figcap{6--71}
\end{bkfigure}}

% ==================================================== Figs 6-72, 6-73, 6-74
% Book p.101.  6-72: the braced gate (a picture, not a lettered diagram) --
%   post at the left, rectangular gate, diagonal brace, ground hatching.
%  6-73: kite A B C D with the diagonal B--D.
%  6-74: parallelogram A B C D with the diagonal A--C.
\newcommand{\FIGVISEVENTYTWOSEVENTYTHREESEVENTYFOUR}{\begin{bkfigure}[\linewidth]
\begin{minipage}{0.48\linewidth}\centering
\begin{tikzpicture}[scale=0.880]
  % post
  \draw[given] (0.10,-0.12) rectangle (0.42,1.92);
  % gate
  \draw[given] (0.52,0.12) rectangle (3.02,1.62);
  \draw[given] (0.66,0.26) rectangle (2.88,1.48);
  % brace
  \draw[given] (0.66,0.26) -- (2.88,1.48);
  % hinges
  \draw[given] (0.42,1.42) -- (0.66,1.42);
  \draw[given] (0.42,0.36) -- (0.66,0.36);
  % ground
  \draw[given] (-0.10,-0.12) -- (3.32,-0.12);
  \foreach \x in {-0.02,0.16,...,3.20}{
    \draw[tick] (\x,-0.12) -- ($(\x,-0.12)+(-0.10,-0.14)$);}
\end{tikzpicture}
\figcap{6--72}
\end{minipage}\hfill
\begin{minipage}{0.22\linewidth}\centering
\begin{tikzpicture}[scale=0.880]
  \coordinate (B) at (0.92,1.72);
  \coordinate (A) at (0,1.02);
  \coordinate (C) at (1.84,1.02);
  \coordinate (D) at (0.92,0);
  \draw[given] (A) -- (B) -- (C) -- (D) -- cycle;
  \draw[given] (B) -- (D);
  \node[vlab, above] at (B) {$B$};
  \node[vlab, left]  at (A) {$A$};
  \node[vlab, right] at (C) {$C$};
  \node[vlab, below] at (D) {$D$};
\end{tikzpicture}
\figcap{6--73}
\end{minipage}\hfill
\begin{minipage}{0.26\linewidth}\centering
\begin{tikzpicture}[scale=0.880]
  \coordinate (D) at (0,0);
  \coordinate (C) at (2.12,0.22);
  \coordinate (B) at (1.82,1.12);
  \coordinate (A) at (0.32,1.02);
  \draw[given] (D) -- (C) -- (B) -- (A) -- cycle;
  \draw[given] (A) -- (C);
  \node[vlab, below left]  at (D) {$D$};
  \node[vlab, right] at (C) {$C$};
  \node[vlab, above right] at (B) {$B$};
  \node[vlab, above left]  at (A) {$A$};
\end{tikzpicture}
\figcap{6--74}
\end{minipage}
\end{bkfigure}}

% =========================================================== Figs 6-75, 6-76
% Book p.102.  6-75: triangle A C with apex B and cevian BD.
%   6-76: the kite A B C D with spine A--C through E, angles 1 and 2 at A.
\newcommand{\FIGVISEVENTYFIVESEVENTYSIX}{\begin{bkfigure}[\linewidth]
\begin{minipage}{0.44\linewidth}\centering
\begin{tikzpicture}[scale=0.938]
  % Exercise 4 gives AB = BC and AD = DC, so D is the midpoint of AC and B
  % stands over it, as the book draws it.
  \coordinate (A) at (0,0);
  \coordinate (C) at (2.62,0);
  \coordinate (B) at (1.31,1.72);
  \coordinate (D) at ($(A)!0.5!(C)$);
  \draw[given] (A) -- (C) -- (B) -- cycle;
  \draw[given] (B) -- (D);
  \node[vlab, left]  at (A) {$A$};
  \node[vlab, right] at (C) {$C$};
  \node[vlab, above] at (B) {$B$};
  \node[vlab, below] at (D) {$D$};
\end{tikzpicture}
\figcap{6--75}
\end{minipage}\hfill
\begin{minipage}{0.52\linewidth}\centering
\begin{tikzpicture}[scale=0.938]
  % Exercise 5 gives AB = AD and BC = CD, so A and C both lie on the
  % perpendicular bisector of BD: B and D must straddle the line AC evenly.
  % They did not -- BD's midpoint sat 0.05 below AC -- and E was placed by a
  % hand-picked ratio instead of being taken as the actual crossing.
  \coordinate (A) at (0,0.86);
  \coordinate (C) at (3.42,0.86);
  \coordinate (B) at (1.22,1.72);
  \coordinate (D) at (1.22,0);
  \draw[given] (A) -- (B) -- (C) -- (D) -- cycle;
  \draw[given] (A) -- (C);
  \draw[given] (B) -- (D);
  \coordinate (E) at (intersection of A--C and B--D);
  \node[vlab, left]  at (A) {$A$};
  \node[vlab, right] at (C) {$C$};
  \node[vlab, above] at (B) {$B$};
  \node[vlab, below] at (D) {$D$};
  \node[vlab, above right] at (E) {$E$};
  \VInumb{C}{A}{B}{0.86cm}{1}
  \VInumb{D}{A}{C}{0.84cm}{2}
\end{tikzpicture}
\figcap{6--76}
\end{minipage}
\end{bkfigure}}

% =========================================================== Figs 6-77, 6-78
% Book p.102.  6-77: parallelogram A B C D with E on BC-side, F interior, and
%   the segments A--C, B--D drawn.  6-78: triangle A C with apex B, E and D on
%   the base.
\newcommand{\FIGVISEVENTYSEVENSEVENTYEIGHT}{\begin{bkfigure}[\linewidth]
\begin{minipage}{0.52\linewidth}\centering
\begin{tikzpicture}[scale=1.179]
  \coordinate (A) at (0,0);
  \coordinate (D) at (2.62,0);
  \coordinate (B) at (0.52,1.12);
  \coordinate (C) at (3.14,1.12);
  \draw[given] (A) -- (D) -- (C) -- (B) -- cycle;
  \draw[given] (B) -- (D);
  \draw[given] (A) -- (C);
  % AF = EC is given, so F and E sit symmetrically on the diagonal.
  \coordinate (F) at ($(A)!0.30!(C)$);
  \coordinate (E) at ($(A)!0.70!(C)$);
  \node[vlab, below left]  at (A) {$A$};
  \node[vlab, below right] at (D) {$D$};
  \node[vlab, above left]  at (B) {$B$};
  \node[vlab, above right] at (C) {$C$};
  % E and F both sit ON the diagonal, so each letter goes on the side away
  % from the long edge it is nearest.
  \node[vlab, below right] at (E) {$E$};
  \node[vlab, above left] at (F) {$F$};
\end{tikzpicture}
\figcap{6--77}
\end{minipage}\hfill
\begin{minipage}{0.44\linewidth}\centering
\begin{tikzpicture}[scale=0.898]
  % Exercise 7: triangle ABC isosceles with vertex B (so B stands over the
  % midpoint of AC) and AD = EC (so E and D are symmetric about it).
  \coordinate (A) at (0,0);
  \coordinate (C) at (2.92,0);
  \coordinate (B) at (1.46,1.72);
  \coordinate (E) at ($(A)!0.32!(C)$);
  \coordinate (D) at ($(A)!0.68!(C)$);
  \draw[given] (A) -- (C) -- (B) -- cycle;
  \draw[given] (B) -- (E);
  \draw[given] (B) -- (D);
  \node[vlab, left]  at (A) {$A$};
  \node[vlab, right] at (C) {$C$};
  \node[vlab, above] at (B) {$B$};
  \node[vlab, below] at (E) {$E$};
  \node[vlab, below] at (D) {$D$};
\end{tikzpicture}
\figcap{6--78}
\end{minipage}
\end{bkfigure}}

% =========================================================== Figs 6-79, 6-80
% Book p.102.  6-79: the re-entrant "dart" A B C D.
%   6-80: triangle A B with apex C, D on the base, E, G, F on the cevian rows,
%   angles 1 and 2 at D.
\newcommand{\FIGVISEVENTYNINEEIGHTY}{\begin{bkfigure}[\linewidth]
\begin{minipage}{0.48\linewidth}\centering
\begin{tikzpicture}[scale=0.938]
  % Exercise 8 gives AB = AD and BC = CD, so A and C lie on the perpendicular
  % bisector of BD -- both at the height of BD's midpoint, 0.81.
  \coordinate (A) at (0,0.81);
  \coordinate (B) at (3.02,1.62);
  \coordinate (D) at (3.02,0);
  \coordinate (C) at (2.12,0.81);
  \draw[given] (A) -- (B) -- (C) -- (D) -- cycle;
  \node[vlab, left]  at (A) {$A$};
  \node[vlab, right] at (B) {$B$};
  \node[vlab, right] at (D) {$D$};
  \node[vlab, left] at (C) {$C$};
\end{tikzpicture}
\figcap{6--79}
\end{minipage}\hfill
\begin{minipage}{0.48\linewidth}\centering
\begin{tikzpicture}[scale=0.938]
  \coordinate (A) at (0,0);
  \coordinate (B) at (3.02,0);
  \coordinate (C) at (1.51,1.72);
  \coordinate (D) at ($(A)!0.5!(B)$);
  \coordinate (E) at ($(A)!0.54!(C)$);
  \coordinate (F) at ($(B)!0.54!(C)$);
  \draw[given] (A) -- (B) -- (C) -- cycle;
  \draw[given] (C) -- (D);
  \draw[given] (E) -- (F);
  % G is where C--D cuts E--F; it was placed as the midpoint of EF instead of
  % being derived, and its letter, set up-and-right, ran into the side CB.
  \coordinate (G) at (intersection of C--D and E--F);
  \node[vlab, left]  at (A) {$A$};
  \node[vlab, right] at (B) {$B$};
  \node[vlab, above] at (C) {$C$};
  \node[vlab, below] at (D) {$D$};
  \node[vlab, above left]  at (E) {$E$};
  \node[vlab, above right] at (F) {$F$};
  \node[vlab] at ($(G)+(-45:0.42)$) {$G$};
  % The book strikes ONE arc across the whole straight angle ADB and writes
  % the two numerals inside it, either side of DC.
  \draw[fig, line width=0.55pt] ([shift=(0:0.36cm)]D)
       arc[start angle=0, end angle=180, radius=0.36cm];
  \VInum{D}{145}{0.92cm}{2}
  \VInum{D}{35}{0.92cm}{1}
\end{tikzpicture}
\figcap{6--80}
\end{minipage}
\end{bkfigure}}

% =========================================================== Figs 6-81, 6-82
% Book p.102.  6-81: triangle A B with D on AB and apex C below, E and F on the
%   line through C; angles 1 and 2 at C.  6-82: segments C--B and D--A crossing.
\newcommand{\FIGVIEIGHTYONEEIGHTYTWO}{\begin{bkfigure}[\linewidth]
\begin{minipage}{0.50\linewidth}\centering
\begin{tikzpicture}[scale=0.958]
  \coordinate (A) at (0,1.42);
  \coordinate (B) at (2.82,1.42);
  \coordinate (C) at (1.41,0);
  \coordinate (D) at ($(A)!0.5!(B)$);
  \draw[given] (A) -- (B);
  \draw[given] (A) -- (C) -- (B);
  \draw[given] (C) -- (D);
  \draw[given] (0.12,0) -- (2.70,0);
  \node[vlab, left]  at (A) {$A$};
  \node[vlab, right] at (B) {$B$};
  \node[vlab, below] at (C) {$C$};
  \node[vlab, above] at (D) {$D$};
  \node[vlab, left]  at (0.12,0) {$E$};
  \node[vlab, right] at (2.70,0) {$F$};
  % One arc across the whole straight angle ECF, numerals inside it either
  % side of CD.  At 0.46 they sat on the line EF itself.
  \draw[fig, line width=0.55pt] ([shift=(0:0.44cm)]C)
       arc[start angle=0, end angle=180, radius=0.44cm];
  \VInum{C}{112}{0.74cm}{1}
  \VInum{C}{68}{0.74cm}{2}
\end{tikzpicture}
\figcap{6--81}
\end{minipage}\hfill
\begin{minipage}{0.46\linewidth}\centering
\begin{tikzpicture}[scale=0.958]
  % Exercise 11 gives AC = BD and AD = BC, so C and D must sit symmetrically
  % over AB; the draft's D was 0.18 off, breaking both.
  \coordinate (A) at (0,0);
  \coordinate (B) at (2.82,0);
  \coordinate (C) at (0.72,1.32);
  \coordinate (D) at (2.10,1.32);
  \draw[given] (A) -- (B);
  \draw[given] (A) -- (C);
  \draw[given] (A) -- (D);
  \draw[given] (C) -- (B);
  \draw[given] (D) -- (B);
  \node[vlab, left]  at (A) {$A$};
  \node[vlab, right] at (B) {$B$};
  \node[vlab, above] at (C) {$C$};
  \node[vlab, above] at (D) {$D$};
\end{tikzpicture}
\figcap{6--82}
\end{minipage}
\end{bkfigure}}

% =========================================================== Figs 6-83, 6-84
% Book p.103.  6-83: A at the top with B, C at the ends and D, E crossing.
%   6-84: triangle A B with apex C, D on the base, O interior.
\newcommand{\FIGVIEIGHTYTHREEEIGHTYFOUR}{\begin{bkfigure}[\linewidth]
\begin{minipage}{0.46\linewidth}\centering
\begin{tikzpicture}[scale=0.987]
  % Exercise 12 gives AB = BC, so A and C must be mirror images across the
  % horizontal through B; the draft had them 0.80 and 1.02 from it.
  \coordinate (A) at (2.02,1.82);
  \coordinate (B) at (0,1.02);
  \coordinate (C) at (2.02,0.22);
  \coordinate (D) at ($(B)!0.52!(A)$);
  \coordinate (E) at ($(B)!0.52!(C)$);
  \draw[given] (B) -- (A);
  \draw[given] (B) -- (C);
  \draw[given] (A) -- (E);
  \draw[given] (C) -- (D);
  \node[vlab, above] at (A) {$A$};
  \node[vlab, left]  at (B) {$B$};
  \node[vlab, below right] at (C) {$C$};
  \node[vlab, above left]  at (D) {$D$};
  \node[vlab, below left]  at (E) {$E$};
\end{tikzpicture}
\figcap{6--83}
\end{minipage}\hfill
\begin{minipage}{0.50\linewidth}\centering
\begin{tikzpicture}[scale=0.987]
  \coordinate (A) at (0,0);
  \coordinate (B) at (3.02,0);
  \coordinate (C) at (1.51,1.72);
  \coordinate (D) at ($(A)!0.5!(B)$);
  \draw[given] (A) -- (B) -- (C) -- cycle;
  \draw[given] (C) -- (D);
  \coordinate (Pb) at ($(B)!0.42!(C)$);
  \coordinate (Pa) at ($(A)!0.42!(C)$);
  \draw[given] (A) -- (Pb);
  \draw[given] (B) -- (Pa);
  \coordinate (O) at (intersection of C--D and A--Pb);
  \node[vlab, left]  at (A) {$A$};
  \node[vlab, right] at (B) {$B$};
  \node[vlab, above] at (C) {$C$};
  \node[vlab, below] at (D) {$D$};
  % Six rays leave O; the only wide gap is the one opening to the right,
  % between the two cevians.  "above right" lay along A--Pb.
  \node[vlab] at ($(O)+(56:0.62)$) {$O$};
\end{tikzpicture}
\figcap{6--84}
\end{minipage}
\end{bkfigure}}

% =========================================================== Figs 6-85, 6-86
% Book p.103.  6-85: triangle A C with B above A, D and E on the cevians,
%   angles 1 and 2 at A/B.  6-86: parallelogram D C A B with E, F on the sides
%   and M the crossing of the two cevians.
\newcommand{\FIGVIEIGHTYFIVEEIGHTYSIX}{\begin{bkfigure}[\linewidth]
\begin{minipage}{0.48\linewidth}\centering
% Drawn at 1.05, not 0.80.  Both numerals hang off the SAME short side AB --
% 1 in the wedge at A, 2 in the wedge at B -- and each wedge is only about 56
% degrees wide, so the numeral has to sit far enough out to clear both its
% rays while staying clear of the other numeral coming the other way.  At 0.80
% there was no such distance: 1 grazed A--E (0.00 pt).  Enlarging the plate
% buys the room without moving either numeral off its bisector.
\begin{tikzpicture}[scale=1.194]
  % Exercise 14 gives AC = BC, so C must sit on the perpendicular bisector of
  % AB.  The draft's C gave AC = 3.08 against BC = 2.91.
  \coordinate (A) at (0,0);
  \coordinate (B) at (0.22,1.42);
  \coordinate (C) at (3.07,0.25);
  \coordinate (D) at ($(A)!0.52!(C)$);
  \coordinate (E) at ($(B)!0.52!(C)$);
  \draw[given] (A) -- (B);
  \draw[given] (A) -- (C);
  \draw[given] (B) -- (C);
  \draw[given] (A) -- (E);
  \draw[given] (B) -- (D);
  \node[vlab, below] at (A) {$A$};
  \node[vlab, above] at (B) {$B$};
  \node[vlab, right] at (C) {$C$};
  \node[vlab, below] at (D) {$D$};
  \node[vlab, above] at (E) {$E$};
  % 1 sits at A (angle EAB) and 2 at B (angle ABD) -- opposite ends of the same
  % short side, each pushed along a bisector that runs INTO the other.  AB is
  % only 1.44 units, and 0.66 + 0.90 overshot it, so the two numerals crossed
  % and printed as "12".  Measured off the plate, the book puts 1 about 0.31 of
  % AB above A and 2 about 0.41 of AB below B, leaving a clear gap between them.
  \VInumb{E}{A}{B}{0.45cm}{1}
  \VInumb{A}{B}{D}{0.59cm}{2}
\end{tikzpicture}
\figcap{6--85}
\end{minipage}\hfill
\begin{minipage}{0.48\linewidth}\centering
\begin{tikzpicture}[scale=0.909]
  \coordinate (A) at (0,0);
  \coordinate (B) at (2.62,0);
  \coordinate (D) at (0.52,1.22);
  \coordinate (C) at (3.14,1.22);
  % AM = MC is given, so M must be the midpoint of the diagonal -- which
  % happens exactly when E and F are symmetric about the parallelogram's
  % centre, i.e. when their two ratios add to 1.  The draft's 0.30 and 0.62
  % did not, so the drawing contradicted its own hypothesis.
  \coordinate (E) at ($(D)!0.38!(C)$);
  \coordinate (F) at ($(A)!0.62!(B)$);
  \draw[given] (A) -- (B) -- (C) -- (D) -- cycle;
  \draw[given] (A) -- (C);
  \draw[given] (E) -- (F);
  \coordinate (M) at (intersection of A--C and E--F);
  \node[vlab, below left]  at (A) {$A$};
  \node[vlab, below right] at (B) {$B$};
  \node[vlab, above left]  at (D) {$D$};
  \node[vlab, above right] at (C) {$C$};
  \node[vlab, above] at (E) {$E$};
  \node[vlab, below] at (F) {$F$};
  \node[vlab, above left] at (M) {$M$};
\end{tikzpicture}
\figcap{6--86}
\end{minipage}
\end{bkfigure}}

% =========================================================== Figs 6-87, 6-88
% Book p.103.  6-87: triangle A B with apex C, D and E on the sides, crossing
%   cevians AE and BD; angles 1..4.  6-88: the kite A B D C with E below.
\newcommand{\FIGVIEIGHTYSEVENEIGHTYEIGHT}{\begin{bkfigure}[\linewidth]
\begin{minipage}{0.52\linewidth}\centering
\begin{tikzpicture}[scale=1.381]
  \coordinate (A) at (0,0);
  \coordinate (B) at (3.32,0);
  \coordinate (C) at (1.66,1.72);
  \coordinate (D) at ($(A)!0.58!(C)$);
  \coordinate (E) at ($(B)!0.58!(C)$);
  \draw[given] (A) -- (B) -- (C) -- cycle;
  \draw[given] (A) -- (E);
  \draw[given] (B) -- (D);
  \node[vlab, left]  at (A) {$A$};
  \node[vlab, right] at (B) {$B$};
  \node[vlab, above] at (C) {$C$};
  \node[vlab, above left]  at (D) {$D$};
  \node[vlab, above right] at (E) {$E$};
  % All four angles sit at the BASE vertices: the cevians split the angle at
  % A into 1 (next to AB) and 3 (next to AC), and the angle at B into 2 and
  % 4.  The draft put 3 at D and 4 at E, which is not what the plate shows
  % and left Exercise 16 marking angles it never names.
  \VIarc{B}{A}{E}{0.66cm}   \VInumb{B}{A}{E}{0.94cm}{1}
  \VIarc{A}{B}{D}{0.66cm}   \VInumb{A}{B}{D}{0.94cm}{2}
  \VIarc{E}{A}{C}{0.40cm}   \VInumb{E}{A}{C}{0.98cm}{3}
  \VIarc{D}{B}{C}{0.40cm}   \VInumb{D}{B}{C}{0.98cm}{4}
\end{tikzpicture}
\figcap{6--87}
\end{minipage}\hfill
\begin{minipage}{0.44\linewidth}\centering
% Proportions remeasured against scan07 p18 (book p.103) at 150 dpi.  The draft
% kite was twice as wide for its height as the book's (height/width 0.96 against
% the plate's 1.76) and put D barely below BC, which is what drove D's letter
% into the numeral 2.  Measured off the plate, with BC as the unit: A sits 0.520
% above BC, D 0.495 below it, E 1.244 below it.
% Everything is mirror-symmetric about the vertical AE, so BE = CE and angle 1 =
% angle 2 (the hypotheses) and angle 3 = angle 4 (the conclusion) all hold
% exactly rather than by eye.
\begin{tikzpicture}[scale=1.062]
  \coordinate (B) at (0,0);
  \coordinate (C) at (2.44,0);
  \coordinate (A) at (1.22,1.269);
  \coordinate (D) at (1.22,-1.208);
  \coordinate (E) at (1.22,-3.035);
  \draw[given] (B) -- (A) -- (C);
  \draw[given] (B) -- (D) -- (C);
  \draw[given] (B) -- (E) -- (C);
  \draw[given] (A) -- (E);
  \node[vlab, above] at (A) {$A$};
  \node[vlab, left]  at (B) {$B$};
  \node[vlab, right] at (C) {$C$};
  % The plate letters D in the wedge between DA and DC, just up and right of the
  % point -- the only quadrant at D that no stroke crosses.  That wedge is only
  % 45 degrees wide (DA runs at 90, DC at 44.7), so "above right" -- the 45
  % degree diagonal -- laid the letter exactly along DC.  Measured on the plate
  % the letter sits at 67.7 degrees, which is the bisector; place it there.
  \node[vlab] at ($(D)+(70:0.58)$) {$D$};
  \node[vlab, below] at (E) {$E$};
  % 3 and 4 are at B and C -- the angles ABD and ACD that Exercise 17 asks
  % about -- not at D, where an earlier draft put them.  1 and 2 stay at E.
  % Distances measured off the plate: numerals at 0.33 of BC from B and C, arcs
  % at 0.21; at E the numerals sit 0.39 of BC out and the arcs at 0.28, so there
  % the numeral is OUTSIDE its arc, as the book draws it.
  \VIarc{A}{B}{D}{0.52cm}   \VInumb{A}{B}{D}{0.80cm}{3}
  \VIarc{A}{C}{D}{0.52cm}   \VInumb{A}{C}{D}{0.80cm}{4}
  \VIarc{B}{E}{A}{0.68cm}   \VInumb{B}{E}{A}{0.96cm}{1}
  \VIarc{A}{E}{C}{0.68cm}   \VInumb{A}{E}{C}{0.96cm}{2}
\end{tikzpicture}
\figcap{6--88}
\end{minipage}
\end{bkfigure}}

% ================================================================= Fig 6-89
% Book p.103.  Parallelogram D C A B with E and F on the diagonal, angles 1, 2.
\newcommand{\FIGVIEIGHTYNINE}{\begin{bkfigure}[\linewidth]
\begin{tikzpicture}[scale=1.118]
  \coordinate (A) at (0,0);
  \coordinate (B) at (2.82,0);
  \coordinate (D) at (0.52,1.22);
  \coordinate (C) at (3.34,1.22);
  % The book draws the diagonal A--C and hangs D--E and B--F off it, with E
  % and F on that diagonal; there is no D--B.  Angle 1 = ADE at D and
  % angle 2 = CBF at B are the pair Exercise 19 gives, and DE = BF is what it
  % asks for -- so both of those segments have to be in the picture.  The
  % draft drew D--B instead and left E and F unjoined to anything.
  \draw[given] (A) -- (B) -- (C) -- (D) -- cycle;
  \draw[given] (A) -- (C);
  \coordinate (E) at ($(A)!0.30!(C)$);
  \coordinate (F) at ($(A)!0.70!(C)$);
  \draw[given] (D) -- (E);
  \draw[given] (B) -- (F);
  \node[vlab, below left]  at (A) {$A$};
  \node[vlab, below right] at (B) {$B$};
  \node[vlab, above left]  at (D) {$D$};
  \node[vlab, above right] at (C) {$C$};
  \node[vlab] at ($(E)+(82:0.46)$) {$E$};
  \node[vlab] at ($(F)+(262:0.46)$) {$F$};
  \VIarc{A}{D}{E}{0.36cm}   \VInumb{A}{D}{E}{0.62cm}{1}
  \VIarc{C}{B}{F}{0.36cm}   \VInumb{C}{B}{F}{0.62cm}{2}
\end{tikzpicture}
\figcap{6--89}
\end{bkfigure}}

% =========================================================== Figs 6-90, 6-91
% Book p.104.  6-90: trapezoid A B C D with both diagonals crossing at E.
%   6-91: parallelogram D C A B with E, F on the diagonal, angles 1..4.
\newcommand{\FIGVININETYNINETYONE}{\begin{bkfigure}[\linewidth]
\begin{minipage}{0.48\linewidth}\centering
\begin{tikzpicture}[scale=0.873]
  % Exercise 20 gives AD = CB and BD = AC, so the trapezoid is isosceles;
  % the draft's A and B were 0.52 and 0.50 in from the ends.
  \coordinate (D) at (0,0);
  \coordinate (C) at (2.92,0);
  \coordinate (A) at (0.51,1.12);
  \coordinate (B) at (2.41,1.12);
  \draw[given] (D) -- (C) -- (B) -- (A) -- cycle;
  \draw[given] (A) -- (C);
  \draw[given] (D) -- (B);
  \coordinate (E) at (intersection of A--C and D--B);
  \node[vlab, below left]  at (D) {$D$};
  \node[vlab, below right] at (C) {$C$};
  \node[vlab, above left]  at (A) {$A$};
  \node[vlab, above right] at (B) {$B$};
  \node[vlab, below] at (E) {$E$};
\end{tikzpicture}
\figcap{6--90}
\end{minipage}\hfill
\begin{minipage}{0.48\linewidth}\centering
\begin{tikzpicture}[scale=1.254]
  \coordinate (A) at (0,0);
  \coordinate (B) at (2.82,0);
  \coordinate (D) at (0.52,1.22);
  \coordinate (C) at (3.34,1.22);
  % Same construction as Fig. 6-89, and checked against the same plate: the
  % diagonal A--C with D--E and B--F hung off it, and NO D--B.  Exercise 21
  % gives AB = CD, AE = CF and angle 3 = angle 4, and asks for angle 1 =
  % angle 2, where 1 = ADE at D, 2 = CBF at B, 3 = BAC at A and 4 = DCA at C.
  % The draft drew D--B, left E and F unjoined, and marked four different
  % angles.
  \draw[given] (A) -- (B) -- (C) -- (D) -- cycle;
  \draw[given] (A) -- (C);
  \coordinate (E) at ($(A)!0.30!(C)$);
  \coordinate (F) at ($(A)!0.70!(C)$);
  \draw[given] (D) -- (E);
  \draw[given] (B) -- (F);
  \node[vlab, below left]  at (A) {$A$};
  \node[vlab, below right] at (B) {$B$};
  \node[vlab, above left]  at (D) {$D$};
  \node[vlab, above right] at (C) {$C$};
  \node[vlab] at ($(E)+(70:0.34)$) {$E$};
  \node[vlab] at ($(F)+(250:0.34)$) {$F$};
  \VIarc{A}{D}{E}{0.36cm}   \VInumb{A}{D}{E}{0.62cm}{1}
  \VIarc{C}{B}{F}{0.36cm}   \VInumb{C}{B}{F}{0.62cm}{2}
  \VIarc{B}{A}{C}{0.36cm}   \VInumb{B}{A}{C}{1.24cm}{3}
  \VIarc{D}{C}{A}{0.36cm}   \VInumb{D}{C}{A}{1.24cm}{4}
\end{tikzpicture}
\figcap{6--91}
\end{minipage}
\end{bkfigure}}

% ================================================================= Fig 6-92
% Book p.104, Exercise 22.  Rebuilt against scan07 p19 (book p.104).
% The book draws the OUTER PENTAGON A B E D C as well as the five chords --
% an earlier draft drew the bare five-pointed star, leaving the outer edges
% A--B, B--E, E--D, D--C, C--A off the page even though the exercise's two
% triangles ABC and ADE are built from them.
% The pentagon is NOT regular: measured off the plate it is markedly wider
% than tall (width 2.84 against height 2.37).  It IS mirror-symmetric about
% the vertical through A, which is what makes AB = AC, AD = AE and
% angle 1 = angle 2 hold exactly -- the hypotheses of Exercise 22.
\newcommand{\FIGVININETYTWO}{\begin{bkfigure}[\linewidth]
\begin{tikzpicture}[scale=1.612]
  \coordinate (A) at (0,0);
  \coordinate (B) at (-1.42,1.214);
  \coordinate (C) at (1.42,1.214);
  \coordinate (E) at (-0.81,2.366);
  \coordinate (D) at (0.81,2.366);
  \draw[given] (A) -- (B) -- (E) -- (D) -- (C) -- cycle;   % outer pentagon
  \draw[given] (B) -- (C);                                 % the five chords
  \draw[given] (B) -- (D);
  \draw[given] (C) -- (E);
  \draw[given] (A) -- (D);
  \draw[given] (A) -- (E);
  % Inner pentagon: every vertex is the crossing of two chords, so each is
  % derived, never placed by hand.
  \coordinate (F) at (intersection of B--D and C--E);
  \coordinate (G) at (intersection of B--D and A--E);
  \coordinate (H) at (intersection of A--D and C--E);
  \coordinate (I) at (intersection of A--E and B--C);
  \coordinate (J) at (intersection of A--D and B--C);
  \node[vlab, above left]  at (E) {$E$};
  \node[vlab, above right] at (D) {$D$};
  \node[vlab, left]  at (B) {$B$};
  \node[vlab, right] at (C) {$C$};
  \node[vlab, below] at (A) {$A$};
  % Label sides re-measured on scan07 p19 (book p.104) at 4x, by solving each
  % crossing from the plate's own vertices rather than eyeballing it:
  %   F point (568,266), letter (565,210) -> directly ABOVE   (draft: below)
  %   G point (369,357), letter (413,383) -> down and RIGHT   (draft: above left)
  %   H point (742,363), letter (680,390) -> down and LEFT    (draft: above right)
  % G and H both lean INWARD, into the wedge under the chord that makes them;
  % the draft had all three on the opposite side of their own point.
  \node[vlab] at ($(F)+(90:0.30)$) {$F$};
  \node[vlab] at ($(G)+(330:0.32)$) {$G$};
  \node[vlab] at ($(H)+(205:0.34)$) {$H$};
  \node[vlab, below left]  at (I) {$I$};
  \node[vlab, below right] at (J) {$J$};
  % Angle 1 = BAE and angle 2 = DAC (the two given angles at the vertex);
  % angle 3 = DBC and angle 4 = ECB (the pair to be proved congruent).
  % Measured on the plate: the numerals at A sit 0.205 of BC out from the
  % vertex, just outside their arcs (0.16 of BC).  The draft put them at 0.35
  % and 0.32 of BC, which carried numeral 1 up into the letter I -- the two
  % then printed as "I_1" and the collision audit merged them into one word
  % and reported nothing.  See the note on stacked labels in ch06-UNCERTAIN.
  \VIarc{B}{A}{E}{0.46cm}   \VInumb{B}{A}{E}{0.78cm}{1}
  \VIarc{D}{A}{C}{0.46cm}   \VInumb{D}{A}{C}{0.78cm}{2}
  % The wedges at B and C are narrow, so a numeral 1.00 out along the
  % bisector still grazed the chords BD and CE; 0.72 keeps it inside.
  \VIarc{D}{B}{C}{0.46cm}   \VInumb{D}{B}{C}{0.72cm}{3}
  \VIarc{E}{C}{B}{0.46cm}   \VInumb{E}{C}{B}{0.72cm}{4}
\end{tikzpicture}
\figcap{6--92}
\end{bkfigure}}

% ================================================================= Fig 6-93
% Book p.104.  Three copies of "a ray OA from a point O on a line l, with B on
% l": the angle AOB acute, right, and obtuse.
\newcommand{\FIGVININETYTHREE}{\begin{bkfigure}[\linewidth]
\begin{tikzpicture}[scale=1.170]
 \foreach \dx/\ax/\ay in {0/0.86/1.22, 2.62/0.32/1.32, 5.24/-0.32/1.28}{
  \begin{scope}[xshift=\dx cm]
   \coordinate (O) at (0.32,0);
   \coordinate (B) at (1.12,0);
   \draw[given] (-0.22,0) -- (1.62,0);
   \draw[given] (O) -- ($(O)+(\ax,\ay)$);
   \dt{O}
   \dt{B}
   \node[vlab, above right] at ($(O)+(\ax,\ay)$) {$A$};
   \node[vlab, below] at (O) {$O$};
   \node[vlab, above right] at (B) {$B$};
   \node[slab, left] at (-0.22,0) {$l$};
  \end{scope}}
\end{tikzpicture}
\figcap{6--93}
\end{bkfigure}}

% ================================================================= Fig 6-94
% Book p.105.  The three cases of the right-angle construction: C on the same
% side of O as B; C = O; C on the opposite side.  A' is the reflection of A in
% l, and AA' meets l at C.  The circular sweep is the book's arc marking the
% two congruent angles at O.
\newcommand{\FIGVININETYFOUR}{\begin{bkfigure}[\linewidth]
\begin{tikzpicture}[scale=1.118]
  % Arcs rebuilt against scan07 p20 (book p.105) at 150 dpi.  The draft drew a
  % full CIRCLE round O in all three panels; the book draws no circle anywhere
  % here.  It strikes one arc from ray OA round to ray OA' on the side B lies
  % on -- the two supplementary angles the proof pairs up -- which is a straight
  % angle in case 2 and a reflex one in case 3, so \VIarc (always minor) cannot
  % express it and \VIarcvia is used instead.  Case 3 carries a SECOND, smaller
  % arc the other way round, through C's side, because there the proof pairs
  % angle AOC with angle A'OC and those now lie opposite B.
  % --- case 1: C between O and B
  \begin{scope}
   \coordinate (O) at (0.32,0);
   % Measured off the scan, OB is 2.25 times OC; the draft had it at 1.86, which
   % is what pushed B's letter up against C's.
   \coordinate (B) at (1.78,0);
   \coordinate (A) at (1.02,1.42);
   \coordinate (Ap) at (1.02,-1.42);
   \coordinate (C) at (intersection of A--Ap and O--B);
   \draw[given] (-0.12,0) -- (1.96,0);
   \draw[given] (O) -- (A);
   \draw[given] (O) -- (Ap);
   \draw[given] (A) -- (Ap);
   \VIarcvia{A}{O}{Ap}{B}{0.60cm}
   \dt{O} \dt{B}
   \node[slab, left] at (-0.12,0) {$l$};
   % Straight up from O is outside the arc's sector (it spans OA round to OA'
   % through B), so the book's own "just above the point" placement is clear.
   \node[vlab, above left, inner sep=0.6pt] at (O) {$O$};
   \node[vlab, above right] at (C) {$C$};
   % Centred, not above-right: set to the right it reached across the panel gap
   % and came within 0.5pt of the NEXT panel's l.  The plate centres it too.
   \node[vlab, above] at (B) {$B$};
   \node[vlab, above left] at (A) {$A$};
   \node[vlab, below left] at (Ap) {$A\VIpr$};
  \end{scope}
  % --- case 2: C = O
  \begin{scope}[xshift=2.72cm]
   \coordinate (O) at (0.42,0);
   \coordinate (B) at (1.52,0);
   \coordinate (A) at (0.42,1.42);
   \coordinate (Ap) at (0.42,-1.42);
   \draw[given] (-0.12,0) -- (1.92,0);
   \draw[given] (A) -- (Ap);
   \VIarcvia{A}{O}{Ap}{B}{0.62cm}
   \dt{O} \dt{B}
   \node[slab, left] at (-0.12,0) {$l$};
   % O and C coincide here, so the book cannot letter the point in place: it
   % sets the name clear of the drawing, up and to the left, and runs a leader
   % arrow down to the point.  The old draft instead anchored this label
   % south-EAST of a point already left of the panel's own left edge, so the
   % text grew leftwards out of the panel and landed on top of Fig. 6-94's
   % first panel -- which is where the "O=CB" and "C" pile-up came from.
   % Kept clear of the B/C letter band the three panels share: at y = 0.72 this
   % label's box still overlapped the neighbouring panel's B by 0.3pt.
   \node[vlab, anchor=east] at ($(O)+(-0.10,1.02)$) {$O = C$};
   \draw[fig, line width=0.5pt, ->] ($(O)+(-0.34,0.78)$) -- ($(O)+(-0.07,0.11)$);
   \node[vlab, above right] at (B) {$B$};
   \node[vlab, above left] at (A) {$A$};
   \node[vlab, below left] at (Ap) {$A\VIpr$};
  \end{scope}
  % --- case 3: C on the far side of O from B
  \begin{scope}[xshift=5.44cm]
   \coordinate (O) at (1.42,0);
   % B must lie OUTSIDE both arcs.  Measured on the plate (book p.105, photo
   % PDF p.119, panel 3 at 5x): OC = 217 px and OB = 270 px, so OB = 1.244*OC,
   % while the outer arc's radius is 218 px -- it passes essentially through C
   % and B sits clear beyond it.  The draft had OB = 0.667*OC, which put B
   % INSIDE the arc and drove its letter onto the arc stroke.
   \coordinate (B) at (2.54,0);
   \coordinate (A) at (0.52,1.42);
   \coordinate (Ap) at (0.52,-1.42);
   \coordinate (C) at (intersection of A--Ap and O--B);
   \draw[given] (0.02,0) -- (2.76,0);
   \draw[given] (O) -- (A);
   \draw[given] (O) -- (Ap);
   \draw[given] (A) -- (Ap);
   % Measured off the scan: the inner arc crosses l at 0.72 of OC from O, the
   % outer at 0.98 of it.
   \VIarcvia{A}{O}{Ap}{C}{0.65cm}
   \VIarcvia{A}{O}{Ap}{B}{0.88cm}
   \dt{O} \dt{B}
   \node[slab, left] at (0.02,0) {$l$};
   \node[vlab, below right, inner sep=0.6pt] at (O) {$O$};
   \node[vlab, above left]  at (C) {$C$};
   \node[vlab, above] at (B) {$B$};
   \node[vlab, above left] at (A) {$A$};
   \node[vlab, below left] at (Ap) {$A\VIpr$};
  \end{scope}
\end{tikzpicture}
\figcap{6--94}
\end{bkfigure}}

% =========================================================== Figs 6-95, 6-96
% Book p.105.  6-95: P at the left with Q, R, S stacked at the right.
%   6-96: the kite B A D C with spine A--C through E.
\newcommand{\FIGVININETYFIVENINETYSIX}{\begin{bkfigure}[\linewidth]
\begin{minipage}{0.50\linewidth}\centering
\begin{tikzpicture}[scale=1.028]
  \coordinate (P) at (0,0.72);
  \coordinate (Q) at (2.62,1.44);
  \coordinate (R) at (2.62,0.72);
  \coordinate (S) at (2.62,0);
  \draw[given] (P) -- (Q) -- (S) -- cycle;
  \draw[given] (P) -- (R);
  \draw[given] (Q) -- (S);
  \node[vlab, left]  at (P) {$P$};
  \node[vlab, right] at (Q) {$Q$};
  \node[vlab, right] at (R) {$R$};
  \node[vlab, right] at (S) {$S$};
\end{tikzpicture}
\figcap{6--95}
\end{minipage}\hfill
\begin{minipage}{0.46\linewidth}\centering
\begin{tikzpicture}[scale=1.028]
  % Off scan07 p20 the book's kite is narrower than the draft's: the
  % half-width is about 0.6 of the upper axis, and the lower axis is about
  % 1.7 times the upper.  E is taken as the actual crossing, not by ratio.
  \coordinate (A) at (1.32,1.82);
  \coordinate (B) at (0.57,0.92);
  \coordinate (D) at (2.07,0.92);
  \coordinate (C) at (1.32,-0.32);
  \draw[given] (B) -- (A) -- (D) -- (C) -- cycle;
  \draw[given] (B) -- (D);
  \draw[given] (A) -- (C);
  \coordinate (E) at (intersection of B--D and A--C);
  \node[vlab, above] at (A) {$A$};
  \node[vlab, left]  at (B) {$B$};
  \node[vlab, right] at (D) {$D$};
  \node[vlab, below] at (C) {$C$};
  \node[vlab] at ($(E)+(250:0.42)$) {$E$};
\end{tikzpicture}
\figcap{6--96}
\end{minipage}
\end{bkfigure}}

% ================================================================= Fig 6-97
% Book p.106.  The comparison picture for "less than": angle B has been moved
% to the position B' with one side along the first side of A, and its second
% side (dashed) falling interior to A.  A and B are drawn side by side.
\newcommand{\FIGVININETYSEVEN}{\begin{bkfigure}[\linewidth]
\begin{tikzpicture}[scale=1.196]
  % angle A, with the moved copy B' inside it
  \coordinate (A) at (2.42,0);
  \draw[given] (0.10,0) -- (A);
  \draw[given] (A) -- ($(A)+(118:1.52)$);
  \draw[aux]   (A) -- ($(A)+(146:1.32)$);
  \node[vlab, below right] at (A) {$A$};
  \node[slab, above] at ($(A)+(146:1.42)$) {$B\VIpr$};
  % angle B, drawn separately to the right
  \coordinate (B) at (5.62,0);
  \draw[given] (2.92,0) -- (B);
  \draw[given] (B) -- ($(B)+(132:1.62)$);
  \node[vlab, below right] at (B) {$B$};
\end{tikzpicture}
\figcap{6--97}
\end{bkfigure}}

% =========================================================== Figs 6-98, 6-99
% Book p.106.  Definition 6-5: (a) second side of B' interior to A  => B < A;
%   (b) second side of B' exterior to A  => B > A.
\newcommand{\FIGVININETYEIGHTNINETYNINE}{\begin{bkfigure}[\linewidth]
\begin{minipage}{0.48\linewidth}\centering
\begin{tikzpicture}[scale=0.885]
  \coordinate (A) at (2.02,0);
  \draw[given] (0.10,0) -- (A);
  \draw[given] (A) -- ($(A)+(112:1.42)$);
  \draw[aux]   (A) -- ($(A)+(142:1.22)$);
  \node[vlab, below right] at (A) {$A$};
  \node[slab, above] at ($(A)+(142:1.32)$) {$B\VIpr$};
  \coordinate (B) at (4.32,0);
  \draw[given] (2.42,0) -- (B);
  \draw[given] (2.42,0) -- ($(2.42,0)+(38:1.62)$);
  \node[vlab, below right] at (B) {$B$};
\end{tikzpicture}
\figcap{6--98}
\end{minipage}\hfill
\begin{minipage}{0.48\linewidth}\centering
\begin{tikzpicture}[scale=0.885]
  \coordinate (A) at (2.02,0);
  \draw[given] (0.10,0) -- (A);
  \draw[given] (A) -- ($(A)+(122:1.42)$);
  \draw[aux]   (A) -- ($(A)+(96:1.52)$);
  \node[vlab, below right] at (A) {$A$};
  \node[slab, left] at ($(A)+(150:0.92)$) {$B\VIpr$};
  \coordinate (B) at (4.42,0);
  \draw[given] (2.52,0) -- (B);
  \draw[given] (2.52,0) -- ($(2.52,0)+(28:1.72)$);
  \node[vlab, below right] at (B) {$B$};
\end{tikzpicture}
\figcap{6--99}
\end{minipage}
\end{bkfigure}}

% ========================================================= Figs 6-100, 6-101
% Book p.107.  6-100 shows B < A (B' interior to A); 6-101 shows the same pair
% with A moved instead, so that A' is exterior to B.
\newcommand{\FIGVIONEHUNDREDONEHUNDREDONE}{\begin{bkfigure}[\linewidth]
\begin{minipage}{0.48\linewidth}\centering
\begin{tikzpicture}[scale=0.878]
  \coordinate (A) at (2.22,0);
  \draw[given] (0.10,0) -- (A);
  \draw[given] (A) -- ($(A)+(118:1.42)$);
  \draw[aux]   (A) -- ($(A)+(150:1.22)$);
  \node[vlab, below right] at (A) {$A$};
  \node[slab, above] at ($(A)+(150:1.34)$) {$B\VIpr$};
  \coordinate (B) at (4.52,0);
  \draw[given] (2.62,0) -- (B);
  \draw[given] (2.62,0) -- ($(2.62,0)+(32:1.62)$);
  \node[vlab, below right] at (B) {$B$};
\end{tikzpicture}
\figcap{6--100}
\end{minipage}\hfill
\begin{minipage}{0.48\linewidth}\centering
\begin{tikzpicture}[scale=0.878]
  \coordinate (A) at (1.92,0);
  \draw[given] (0.10,0) -- (A);
  \draw[given] (A) -- ($(A)+(126:1.42)$);
  \node[vlab, below right] at (A) {$A$};
  \coordinate (B) at (4.42,0);
  \draw[given] (2.32,0) -- (B);
  \draw[given] (B) -- ($(B)+(152:1.62)$);
  \draw[aux]   (B) -- ($(B)+(122:1.52)$);
  \node[vlab, below right] at (B) {$B$};
  \node[slab, above] at ($(B)+(122:1.62)$) {$A\VIpr$};
\end{tikzpicture}
\figcap{6--101}
\end{minipage}
\end{bkfigure}}

% ================================================================ Fig 6-102
% Book p.107.  The Theorem 6-8 proof: on angle A, points D and E on the sides
% with DE cutting the second side of B' at F; on angle A', the matching D', E'
% with F' chosen so that DF = D'F'.  F is DERIVED as the crossing of DE with
% the second side of B'.
\newcommand{\FIGVIONEHUNDREDTWO}{\begin{bkfigure}[\linewidth]
\begin{tikzpicture}[scale=1.170]
  % --- left: angle A with B' interior
  \coordinate (A) at (2.42,0);
  \coordinate (E) at ($(A)+(150:1.62)$);
  \coordinate (D) at ($(A)+(212:1.62)$);
  \draw[given] (A) -- ($(A)+(150:2.02)$);
  \draw[given] (A) -- ($(A)+(212:2.02)$);
  \draw[given] (D) -- (E);
  \coordinate (Bs) at ($(A)+(174:2.02)$);
  \draw[given] (A) -- (Bs);
  \coordinate (F) at (intersection of D--E and A--Bs);
  \node[vlab, below right] at (A) {$A$};
  \node[vlab, above right] at (E) {$E$};
  % D sits on the lower side of the angle, which runs on past it to
  % 212:2.02 -- so the letter goes square off that ray, not along it.
  \node[vlab] at ($(D)+(302:0.34)$) {$D$};
  \node[vlab, above left] at (F) {$F$};
  \node[slab, above right] at ($(A)+(120:1.13)$) {$B\VIpr$};
  % --- right: angle A' with the matching construction
  \begin{scope}[xshift=3.62cm]
   \coordinate (B) at (2.42,0);
   \coordinate (Ep) at ($(B)+(150:1.62)$);
   \coordinate (Dp) at ($(B)+(212:1.62)$);
   \draw[given] (B) -- ($(B)+(150:2.02)$);
   \draw[given] (B) -- ($(B)+(212:2.02)$);
   \draw[aux]   (Dp) -- (Ep);
   \coordinate (As) at ($(B)+(174:2.02)$);
   \draw[given] (B) -- (As);
   \coordinate (Fp) at (intersection of Dp--Ep and B--As);
   \node[vlab, below right] at (B) {$B$};
   \node[vlab, above right] at (Ep) {$E\VIpr$};
   \node[vlab] at ($(Dp)+(302:0.34)$) {$D\VIpr$};
   \node[vlab, above left] at (Fp) {$F\VIpr$};
   \node[slab, above right] at ($(B)+(132:1.22)$) {$A\VIpr$};
  \end{scope}
\end{tikzpicture}
\figcap{6--102}
\end{bkfigure}}

% ================================================================ Fig 6-103
% Book p.109.  Theorem 6-11: (a) right angle A with supplement A'; (b) the same
% for B and B'; (c) the comparison figure with B'' congruent to B, and B''' its
% supplement, all at one vertex.
\newcommand{\FIGVIONEHUNDREDTHREE}{\begin{bkfigure}[\linewidth]
\begin{tikzpicture}[scale=1.118]
  % (a)
  \begin{scope}
   \coordinate (O) at (0.86,0);
   \draw[given] (0.02,0) -- (1.72,0);
   \draw[given] (O) -- (0.86,1.32);
   \draw[tick] ($(O)+(180:0.52)$) arc[start angle=180, end angle=0, radius=0.52];
   \node[slab] at ($(O)+(136:0.76)$) {$A$};
   \node[slab] at ($(O)+(36:0.81)$)  {$A\VIpr$};
   \node[slab, below] at (0.86,-0.28) {(a)};
  \end{scope}
  % (b)
  \begin{scope}[xshift=2.32cm]
   \coordinate (O) at (0.86,0);
   \draw[given] (0.02,0) -- (1.72,0);
   \draw[given] (O) -- (0.86,1.32);
   \draw[tick] ($(O)+(180:0.52)$) arc[start angle=180, end angle=0, radius=0.52];
   \node[slab] at ($(O)+(136:0.76)$) {$B$};
   \node[slab] at ($(O)+(32:0.80)$)  {$B\VIpr$};
   \node[slab, below] at (0.86,-0.28) {(b)};
  \end{scope}
  % (c)
  \begin{scope}[xshift=4.64cm]
   \coordinate (O) at (1.12,0);
   \draw[given] (0.02,0) -- (2.22,0);
   \draw[given] (O) -- ($(O)+(90:1.32)$);
   \draw[given] (O) -- ($(O)+(62:1.32)$);
   \draw[tick] ($(O)+(180:0.72)$) arc[start angle=180, end angle=0, radius=0.72];
   \node[slab] at ($(O)+(142:1.00)$) {$B\VIppr$};
   \node[slab] at ($(O)+(80:1.24)$) {$A$};
   \node[slab] at ($(O)+(0:1.32)$)  {$A\VIpr$};
   \node[slab] at ($(O)+(30:1.23)$)  {$B\VIpppr$};
   \node[slab, below] at (1.12,-0.28) {(c)};
  \end{scope}
\end{tikzpicture}
\figcap{6--103}
\end{bkfigure}}

% ================================================================ Fig 6-104
% Book p.109.  (a) l perpendicular to m; (b) l the perpendicular bisector of AB.
\newcommand{\FIGVIONEHUNDREDFOUR}{\begin{bkfigure}[\linewidth]
\begin{tikzpicture}[scale=1.170]
  % (a)
  \begin{scope}
   \coordinate (O) at (0.92,0);
   \coordinate (R) at (1.82,0);
   \coordinate (U) at (0.92,1.32);
   \draw[given] (0.02,0) -- (R);
   \draw[given] (O) -- (U);
   \VIrt{R}{O}{U}{6pt}
   \node[slab, above left] at (U) {$l$};
   \node[slab, right] at (R) {$m$};
   \node[slab, below] at (0.92,-0.36) {(a)};
  \end{scope}
  % (b)
  \begin{scope}[xshift=2.92cm]
   \coordinate (A) at (0.02,0);
   \coordinate (B) at (1.92,0);
   \coordinate (M) at ($(A)!0.5!(B)$);
   \coordinate (U) at ($(M)+(0,1.32)$);
   \draw[given] (A) -- (B);
   \draw[given] (M) -- (U);
   \VIrt{B}{M}{U}{6pt}
   \node[slab, above left] at (U) {$l$};
   \node[vlab, left]  at (A) {$A$};
   \node[vlab, right] at (B) {$B$};
   \node[slab, below] at (0.96,-0.36) {(b)};
  \end{scope}
\end{tikzpicture}
\figcap{6--104}
\end{bkfigure}}

% ======================================================= Figs 6-105, 6-106
% Book p.110.  6-105: rectangle A B C D with both diagonals.
%   6-106: triangle A C with apex B and cevian BD.
\newcommand{\FIGVIONEHUNDREDFIVEONEHUNDREDSIX}{\begin{bkfigure}[\linewidth]
\begin{minipage}{0.52\linewidth}\centering
\begin{tikzpicture}[scale=0.955]
  \coordinate (D) at (0,0);
  \coordinate (C) at (2.92,0);
  \coordinate (A) at (0,1.22);
  \coordinate (B) at (2.92,1.22);
  \draw[given] (A) -- (B) -- (C) -- (D) -- cycle;
  \draw[given] (A) -- (C);
  \draw[given] (B) -- (D);
  \node[vlab, above left]  at (A) {$A$};
  \node[vlab, above right] at (B) {$B$};
  \node[vlab, below left]  at (D) {$D$};
  \node[vlab, below right] at (C) {$C$};
\end{tikzpicture}
\figcap{6--105}
\end{minipage}\hfill
\begin{minipage}{0.44\linewidth}\centering
\begin{tikzpicture}[scale=0.955]
  \coordinate (A) at (0,0);
  \coordinate (C) at (2.72,0);
  \coordinate (B) at (1.02,1.62);
  \coordinate (D) at ($(A)!0.44!(C)$);
  \draw[given] (A) -- (C) -- (B) -- cycle;
  \draw[given] (B) -- (D);
  \node[vlab, left]  at (A) {$A$};
  \node[vlab, right] at (C) {$C$};
  \node[vlab, above] at (B) {$B$};
  \node[vlab, below] at (D) {$D$};
\end{tikzpicture}
\figcap{6--106}
\end{minipage}
\end{bkfigure}}

% ======================================================= Figs 6-107, 6-108
% Book p.110.  6-107: rectangle A B C D with E on the base and the two
%   segments AE, BE.  6-108: triangle A C with apex B and cevian BD.
\newcommand{\FIGVIONEHUNDREDSEVENONEHUNDREDEIGHT}{\begin{bkfigure}[\linewidth]
\begin{minipage}{0.52\linewidth}\centering
\begin{tikzpicture}[scale=0.947]
  \coordinate (D) at (0,0);
  \coordinate (C) at (3.02,0);
  \coordinate (A) at (0,1.12);
  \coordinate (B) at (3.02,1.12);
  \coordinate (E) at ($(D)!0.5!(C)$);
  \draw[given] (A) -- (B) -- (C) -- (D) -- cycle;
  \draw[given] (A) -- (E);
  \draw[given] (B) -- (E);
  \node[vlab, above left]  at (A) {$A$};
  \node[vlab, above right] at (B) {$B$};
  \node[vlab, below left]  at (D) {$D$};
  \node[vlab, below right] at (C) {$C$};
  \node[vlab, below] at (E) {$E$};
\end{tikzpicture}
\figcap{6--107}
\end{minipage}\hfill
\begin{minipage}{0.44\linewidth}\centering
\begin{tikzpicture}[scale=0.947]
  \coordinate (A) at (0,0);
  \coordinate (C) at (2.82,0);
  \coordinate (B) at (1.22,1.62);
  \coordinate (D) at ($(A)!0.5!(C)$);
  \draw[given] (A) -- (C) -- (B) -- cycle;
  \draw[given] (B) -- (D);
  \node[vlab, left]  at (A) {$A$};
  \node[vlab, right] at (C) {$C$};
  \node[vlab, above] at (B) {$B$};
  \node[vlab, below] at (D) {$D$};
\end{tikzpicture}
\figcap{6--108}
\end{minipage}
\end{bkfigure}}

% ======================================================= Figs 6-109, 6-110
% Book p.110.  6-109: B--D and A--F crossing at C, with E above.
%   6-110: right triangle D A C with B on DC and E on AD.
\newcommand{\FIGVIONEHUNDREDNINEONEHUNDREDTEN}{\begin{bkfigure}[\linewidth]
\begin{minipage}{0.50\linewidth}\centering
\begin{tikzpicture}[scale=0.960]
  \coordinate (B) at (0.32,1.02);
  \coordinate (D) at (2.62,0.62);
  \coordinate (A) at (0,0);
  \coordinate (E) at (2.92,1.42);
  \draw[given] (B) -- (D);
  \draw[given] (A) -- (E);
  \coordinate (C) at (intersection of B--D and A--E);
  \node[vlab, left]  at (B) {$B$};
  \node[vlab, right] at (D) {$D$};
  \node[vlab, left]  at (A) {$A$};
  \node[vlab, right] at (E) {$E$};
  \node[vlab, below] at (C) {$C$};
\end{tikzpicture}
\figcap{6--109}
\end{minipage}\hfill
\begin{minipage}{0.46\linewidth}\centering
\begin{tikzpicture}[scale=0.960]
  \coordinate (A) at (0,0);
  \coordinate (C) at (2.62,0);
  \coordinate (D) at (0.42,1.72);
  \coordinate (B) at ($(D)!0.52!(C)$);
  \coordinate (E) at ($(A)!0.62!(D)$);
  \draw[given] (A) -- (C) -- (D) -- cycle;
  \draw[given] (A) -- (B);
  \draw[given] (E) -- (B);
  \VIrt{D}{B}{A}{5pt}
  \node[vlab, below left]  at (A) {$A$};
  \node[vlab, right] at (C) {$C$};
  \node[vlab, above] at (D) {$D$};
  \node[vlab, above right] at (B) {$B$};
  \node[vlab, left]  at (E) {$E$};
\end{tikzpicture}
\figcap{6--110}
\end{minipage}
\end{bkfigure}}

% ================================================================ Fig 6-111
% Book p.111.  A, C, D, B with the right angles at C (AC perp CB) and at D
% (AD perp DB); the two triangles share the hypotenuse-like segment A--B.
\newcommand{\FIGVIONEHUNDREDELEVEN}{\begin{bkfigure}[\linewidth]
\begin{tikzpicture}[scale=1.196]
  % Exercise 11 gives AC = AD with AC perp CB and AD perp DB.  Two right
  % angles on AB put C and D on the circle with diameter AB; AC = AD then
  % makes them a mirror pair across AB.  These coordinates are that pair
  % (circle centre (1.47,0.71), radius 1.03276, parameter 100 deg), not an
  % eyeballed quadrilateral.
  \coordinate (A) at (0.72,1.42);
  \coordinate (D) at (2.0390,1.5719);
  \coordinate (C) at (0.6405,0.0947);
  \coordinate (B) at (2.22,0);
  \draw[given] (A) -- (C) -- (B);
  \draw[given] (A) -- (D) -- (B);
  \draw[given] (A) -- (B);
  \VIrt{A}{C}{B}{5pt}
  \VIrt{B}{D}{A}{5pt}
  \VItk{A}{C}{1}
  \VItk{A}{D}{1}
  \node[vlab, above left]  at (A) {$A$};
  \node[vlab, above right] at (D) {$D$};
  \node[vlab, left]  at (C) {$C$};
  \node[vlab, right] at (B) {$B$};
\end{tikzpicture}
\figcap{6--111}
\end{bkfigure}}

% ======================================================= Figs 6-112, 6-113
% Book p.112.  6-112: the two angle-fans of Review exercise 3 -- R, S, T at one
%   vertex and A_1, A_2 at the other, each with a sweeping arc.
%  6-113: two parallels cut by two transversals, points A, B, C, D.
\newcommand{\FIGVIONEHUNDREDTWELVEONEHUNDREDTHIRTEEN}{\begin{bkfigure}[\linewidth]
\begin{minipage}{0.48\linewidth}\centering
% Rebuilt against scan08 p5 (book p.112).  Two corrections to an earlier
% draft, both settled by Review Exercise 3 ("if angle A_1 = angle R and
% angle A_2 = angle S, prove angle B = angle T"), which makes R, S the two
% PARTS of T and A_1, A_2 the two parts of B:
%   (a) the left fan is drawn VERTEX-UP in the book, its three rays running
%       downward; the draft had it vertex-down like the right one.
%   (b) every one of the six names is an ANGLE, marked by its own arc -- R
%       and S on small arcs, T on a large one enclosing both, and likewise
%       A_1, A_2 inside B.  The draft lettered B as a POINT at the tip of the
%       middle ray, which is not what the plate shows.
% Each numeral sits on its own angle's BISECTOR, just outside that angle's
% arc, which is where the book puts them.
\begin{tikzpicture}[scale=0.895]
  % ---- left fan, vertex UP: rays at 238 / 267 / 306 degrees
  \coordinate (V) at (0,0);
  \draw[given] (V) -- ++(238:1.95);
  \draw[given] (V) -- ++(267:1.95);
  \draw[given] (V) -- ++(306:1.95);
  \draw[tick] ([shift=(238:0.66)]V) arc[start angle=238, end angle=267, radius=0.66];
  \draw[tick] ([shift=(267:0.58)]V) arc[start angle=267, end angle=306, radius=0.58];
  \draw[tick] ([shift=(238:1.30)]V) arc[start angle=238, end angle=306, radius=1.30];
  \node[slab] at ($(V)+(228.5:1.52)$) {$R$};
  \node[slab] at ($(V)+(286.5:0.96)$) {$S$};
  \node[slab] at ($(V)+(288:1.62)$)   {$T$};
  % ---- right fan, vertex DOWN: rays at 124 / 87 / 54 degrees
  \begin{scope}[xshift=3.20cm, yshift=-1.95cm]
   \coordinate (W) at (0,0);
   \draw[given] (W) -- ++(124:1.95);
   \draw[given] (W) -- ++(87:1.95);
   \draw[given] (W) -- ++(54:1.95);
   \draw[tick] ([shift=(87:0.58)]W)  arc[start angle=87,  end angle=124, radius=0.58];
   \draw[tick] ([shift=(54:0.66)]W)  arc[start angle=54,  end angle=87,  radius=0.66];
   \draw[tick] ([shift=(54:1.30)]W)  arc[start angle=54,  end angle=124, radius=1.30];
   \node[slab] at ($(W)+(77.5:1.59)$) {$A_1$};
   \node[slab] at ($(W)+(42.5:1.61)$)  {$A_2$};
   \node[slab] at ($(W)+(105.5:1.60)$) {$B$};
  \end{scope}
\end{tikzpicture}
\figcap{6--112}
\end{minipage}\hfill
\begin{minipage}{0.48\linewidth}\centering
\begin{tikzpicture}[scale=0.895]
  \draw[given] (0.02,1.12) -- (3.22,1.12);
  \draw[given] (0.02,0.42) -- (3.22,0.42);
  \coordinate (P) at (0.62,-0.12);
  \coordinate (Q) at (2.42,1.72);
  \draw[given] (P) -- (Q);
  \coordinate (A) at (intersection of P--Q and 0.02,0.42--3.22,0.42);
  \coordinate (B) at (intersection of P--Q and 0.02,1.12--3.22,1.12);
  \coordinate (C) at ($(B)+(0.42,0.16)$);
  \coordinate (D) at ($(A)+(-0.30,-0.42)$);
  \draw[given] (B) -- (C);
  \node[vlab, above left]  at (B) {$B$};
  \node[vlab, above left] at (C) {$C$};
  \node[vlab, below right] at (A) {$A$};
  \node[vlab, below] at (D) {$D$};
\end{tikzpicture}
\figcap{6--113}
\end{minipage}
\end{bkfigure}}

% ======================================================= Figs 6-114, 6-115
% Book p.112.  6-114: A O C on a line with OB perpendicular.
%   6-115: triangles E F G and R S T.
\newcommand{\FIGVIONEHUNDREDFOURTEENONEHUNDREDFIFTEEN}{\begin{bkfigure}[\linewidth]
\begin{minipage}{0.46\linewidth}\centering
\begin{tikzpicture}[scale=0.945]
  \coordinate (A) at (0,0);
  \coordinate (C) at (2.62,0);
  \coordinate (O) at ($(A)!0.5!(C)$);
  \coordinate (B) at ($(O)+(0,1.52)$);
  \draw[given] (A) -- (C);
  \draw[given] (O) -- (B);
  \node[vlab, left]  at (A) {$A$};
  \node[vlab, right] at (C) {$C$};
  \node[vlab, below] at (O) {$O$};
  \node[vlab, above] at (B) {$B$};
\end{tikzpicture}
\figcap{6--114}
\end{minipage}\hfill
\begin{minipage}{0.50\linewidth}\centering
\begin{tikzpicture}[scale=0.945]
  \coordinate (F) at (0,0);
  \coordinate (G) at (1.62,0);
  \coordinate (E) at (0.42,1.22);
  \draw[given] (F) -- (G) -- (E) -- cycle;
  \node[vlab, below left]  at (F) {$F$};
  \node[vlab, below right] at (G) {$G$};
  \node[vlab, above] at (E) {$E$};
  \coordinate (S) at (1.82,0);
  \coordinate (T) at (3.32,0.62);
  \coordinate (R) at (2.32,1.42);
  \draw[given] (S) -- (T) -- (R) -- cycle;
  \node[vlab, below left]  at (S) {$S$};
  \node[vlab, right] at (T) {$T$};
  \node[vlab, above] at (R) {$R$};
\end{tikzpicture}
\figcap{6--115}
\end{minipage}
\end{bkfigure}}

% ======================================================= Figs 6-116, 6-117
% Book p.112.  6-116: triangle R S T.
%  6-117: the river-width figure -- P and Q on opposite banks (the banks drawn
%  as the book's wavy lines), R off the line PQ, with U, S on the sight lines.
\newcommand{\FIGVIONEHUNDREDSIXTEENONEHUNDREDSEVENTEEN}{\begin{bkfigure}[\linewidth]
\begin{minipage}{0.44\linewidth}\centering
\begin{tikzpicture}[scale=0.962]
  \coordinate (T) at (0,0);
  \coordinate (S) at (2.72,0.32);
  \coordinate (R) at (0.72,1.32);
  \draw[given] (T) -- (S) -- (R) -- cycle;
  \node[vlab, below left]  at (T) {$T$};
  \node[vlab, right] at (S) {$S$};
  \node[vlab, above] at (R) {$R$};
\end{tikzpicture}
\figcap{6--116}
\end{minipage}\hfill
\begin{minipage}{0.52\linewidth}\centering
\begin{tikzpicture}[scale=0.962]
  \coordinate (R) at (1.42,1.72);
  \coordinate (U) at (0.32,-0.42);
  \coordinate (T) at (0.02,0);
  \coordinate (S) at (0.72,0);
  \coordinate (P) at (1.72,0.42);
  \coordinate (Q) at (2.92,0.12);
  \draw[given] (R) -- (U);
  \draw[given] (T) -- (P);
  \draw[given] (R) -- (P);
  \draw[aux]   (R) -- (Q);
  \draw[given] (P) -- (Q);
  % the river bank, drawn as the book's wavy rule
  \draw[tick] (2.62,0.92) .. controls (2.82,0.72) and (2.62,0.52) ..
              (2.82,0.32) .. controls (3.02,0.12) and (2.82,-0.08) .. (3.02,-0.28);
  \node[vlab, above] at (R) {$R$};
  \node[vlab, below] at (U) {$U$};
  \node[vlab, left]  at (T) {$T$};
  \node[vlab, below] at (S) {$S$};
  \node[vlab, below] at (P) {$P$};
  \node[vlab, right] at (Q) {$Q$};
\end{tikzpicture}
\figcap{6--117}
\end{minipage}
\end{bkfigure}}

% ======================================================= Figs 6-118, 6-119
% Book p.113.  6-118: triangle A C with apex B and D interior joined to all
%   three vertices.  6-119: triangle A C with apex B, D and E on the base.
\newcommand{\FIGVIONEHUNDREDEIGHTEENONEHUNDREDNINETEEN}{\begin{bkfigure}[\linewidth]
\begin{minipage}{0.46\linewidth}\centering
\begin{tikzpicture}[scale=0.951]
  \coordinate (A) at (0,0);
  \coordinate (C) at (2.62,0);
  \coordinate (B) at (1.31,1.72);
  \coordinate (D) at (1.31,0.62);
  \draw[given] (A) -- (C) -- (B) -- cycle;
  \draw[given] (A) -- (D) -- (C);
  \draw[given] (B) -- (D);
  \node[vlab, left]  at (A) {$A$};
  \node[vlab, right] at (C) {$C$};
  \node[vlab, above] at (B) {$B$};
  % D sits low, so its letter has to share the gap with the base; the book
  % crowds it the same way, and a tighter clearance is what keeps it off AC.
  \node[vlab, below, outer sep=2.0pt] at (D) {$D$};
\end{tikzpicture}
\figcap{6--118}
\end{minipage}\hfill
\begin{minipage}{0.50\linewidth}\centering
\begin{tikzpicture}[scale=0.951]
  \coordinate (A) at (0,0);
  \coordinate (C) at (3.12,0);
  % Exercise 13: ABC isosceles and angle ABD = angle EBC, so D and E are
  % symmetric about the foot of the altitude.
  \coordinate (B) at (1.56,1.62);
  \coordinate (D) at ($(A)!0.32!(C)$);
  \coordinate (E) at ($(A)!0.68!(C)$);
  \draw[given] (A) -- (C) -- (B) -- cycle;
  \draw[given] (B) -- (D);
  \draw[given] (B) -- (E);
  \node[vlab, left]  at (A) {$A$};
  \node[vlab, right] at (C) {$C$};
  \node[vlab, above] at (B) {$B$};
  \node[vlab, below] at (D) {$D$};
  \node[vlab, below] at (E) {$E$};
\end{tikzpicture}
\figcap{6--119}
\end{minipage}
\end{bkfigure}}

% ======================================================= Figs 6-120, 6-121
% Book p.113.  6-120: triangle A C with apex B, D on the base, E and F the
%   midpoints of the sides, forming the medial triangle.
%  6-121: Q S R on a line with PS perpendicular at S.
\newcommand{\FIGVIONEHUNDREDTWENTYONEHUNDREDTWENTYONE}{\begin{bkfigure}[\linewidth]
\begin{minipage}{0.52\linewidth}\centering
\begin{tikzpicture}[scale=0.964]
  \coordinate (A) at (0,0);
  \coordinate (C) at (3.02,0);
  \coordinate (B) at (1.51,1.72);
  \coordinate (D) at ($(A)!0.5!(C)$);
  \coordinate (E) at ($(A)!0.5!(B)$);
  \coordinate (F) at ($(B)!0.5!(C)$);
  \draw[given] (A) -- (C) -- (B) -- cycle;
  \draw[given] (E) -- (D) -- (F);
  \draw[given] (E) -- (F);
  \node[vlab, left]  at (A) {$A$};
  \node[vlab, right] at (C) {$C$};
  \node[vlab, above] at (B) {$B$};
  \node[vlab, below] at (D) {$D$};
  \node[vlab, above left]  at (E) {$E$};
  \node[vlab, above right] at (F) {$F$};
\end{tikzpicture}
\figcap{6--120}
\end{minipage}\hfill
\begin{minipage}{0.44\linewidth}\centering
\begin{tikzpicture}[scale=0.964]
  \coordinate (Q) at (0,0);
  \coordinate (R) at (2.42,0);
  \coordinate (S) at ($(Q)!0.5!(R)$);
  \coordinate (P) at ($(S)+(0,1.62)$);
  \draw[given] (Q) -- (R);
  \draw[given] (S) -- (P);
  \node[vlab, left]  at (Q) {$Q$};
  \node[vlab, right] at (R) {$R$};
  \node[vlab, below] at (S) {$S$};
  \node[vlab, above] at (P) {$P$};
\end{tikzpicture}
\figcap{6--121}
\end{minipage}
\end{bkfigure}}

% ======================================================= Figs 6-122, 6-123
% Book p.113.  6-122: A--B and C--D crossing at O (O bisects both).
%   6-123: A with B, C, D, E -- the two isosceles triangles of exercise 19.
\newcommand{\FIGVIONEHUNDREDTWENTYTWOONEHUNDREDTWENTYTHREE}{\begin{bkfigure}[\linewidth]
\begin{minipage}{0.50\linewidth}\centering
\begin{tikzpicture}[scale=0.958]
  \coordinate (A) at (0,1.22);
  \coordinate (B) at (3.02,0);
  % O bisects BOTH segments, so C--D's midpoint must land exactly on A--B's;
  % D was 0.02 out.
  \coordinate (C) at (1.92,1.22);
  \coordinate (D) at (1.10,0);
  \coordinate (O) at (intersection of A--B and C--D);
  \draw[given] (A) -- (B);
  \draw[given] (C) -- (D);
  \draw[given] (A) -- (D);
  \draw[given] (C) -- (B);
  \node[vlab, left]  at (A) {$A$};
  \node[vlab, right] at (B) {$B$};
  \node[vlab, above] at (C) {$C$};
  \node[vlab, below] at (D) {$D$};
  \node[vlab, below right] at (O) {$O$};
\end{tikzpicture}
\figcap{6--122}
\end{minipage}\hfill
\begin{minipage}{0.46\linewidth}\centering
\begin{tikzpicture}[scale=0.958]
  % Exercise 19 gives AB = AD, BC = DE and AC = AE.  Those three hold exactly
  % when the fan is struck with |AB| = |AD| = 1.90 and |AC| = |AE| = 1.35 and
  % the two opening angles BAC and DAE are equal (40 degrees each): the rays
  % run at -5, -45, -85 and -125 degrees from A.  The draft's five vertices
  % were eyeballed and satisfied none of the three.
  \coordinate (A) at (0.80,1.95);
  \coordinate (B) at ($(A)+(-5:1.90)$);
  \coordinate (C) at ($(A)+(-45:1.35)$);
  \coordinate (D) at ($(A)+(-85:1.90)$);
  \coordinate (E) at ($(A)+(-125:1.35)$);
  \draw[given] (A) -- (B) -- (C) -- (D) -- (E) -- cycle;
  \draw[given] (A) -- (C);
  \draw[given] (A) -- (D);
  \node[vlab, above] at (A) {$A$};
  \node[vlab, right] at (B) {$B$};
  \node[vlab, below right] at (C) {$C$};
  \node[vlab, below right] at (D) {$D$};
  \node[vlab, left]  at (E) {$E$};
\end{tikzpicture}
\figcap{6--123}
\end{minipage}
\end{bkfigure}}

% ================================================================ Fig 6-124
% Book p.113.  Equilateral triangle A with base C B, D, E, F the midpoints of
% the sides -- the medial triangle EFD of exercise 20.
\newcommand{\FIGVIONEHUNDREDTWENTYFOUR}{\begin{bkfigure}[\linewidth]
\begin{tikzpicture}[scale=1.196]
  \coordinate (C) at (0,0);
  \coordinate (B) at (2.82,0);
  \coordinate (A) at ($(C)!0.5!(B)+(0,2.44)$);
  \coordinate (D) at ($(C)!0.5!(B)$);
  \coordinate (E) at ($(C)!0.5!(A)$);
  \coordinate (F) at ($(A)!0.5!(B)$);
  \draw[given] (C) -- (B) -- (A) -- cycle;
  \draw[given] (E) -- (F);
  \draw[given] (E) -- (D);
  \draw[given] (F) -- (D);
  \node[vlab, left]  at (C) {$C$};
  \node[vlab, right] at (B) {$B$};
  \node[vlab, above] at (A) {$A$};
  \node[vlab, below] at (D) {$D$};
  \node[vlab, left]  at (E) {$E$};
  \node[vlab, right] at (F) {$F$};
\end{tikzpicture}
\figcap{6--124}
\end{bkfigure}}

% ======================================================= Figs 6-125, 6-126
% Book p.114.  6-125: A D with B, C above and the crossing segments AC, BD.
%  6-126: the cevian web -- triangle A B with apex C; D, E on the sides;
%  J, I, K on the base; F, G, H the interior crossings; angles 1, 2.
\newcommand{\FIGVIONEHUNDREDTWENTYFIVEONEHUNDREDTWENTYSIX}{\begin{bkfigure}[\linewidth]
\begin{minipage}{0.46\linewidth}\centering
\begin{tikzpicture}[scale=0.833]
  \coordinate (A) at (0,0);
  \coordinate (D) at (3.02,0);
  \coordinate (B) at (0.72,1.32);
  \coordinate (C) at (2.02,1.32);
  \draw[given] (A) -- (D);
  \draw[given] (A) -- (B) -- (D);
  \draw[given] (A) -- (C) -- (D);
  \node[vlab, left]  at (A) {$A$};
  \node[vlab, right] at (D) {$D$};
  \node[vlab, above] at (B) {$B$};
  \node[vlab, above] at (C) {$C$};
\end{tikzpicture}
\figcap{6--125}
\end{minipage}\hfill
\begin{minipage}{0.50\linewidth}\centering
% Rebuilt against scan08 p8 (book p.114).  The book joins D--J and E--K, not
% D--E.  That matters twice over: the draft's D--E left the exercise's
% segments GJ and HK undrawn, and it made G and H DEGENERATE -- they were
% asked for as (A--E cap D--E) and (B--D cap D--E), two pairs of lines that
% already share an endpoint, so G landed exactly on E and H exactly on D and
% the two letters printed on top of the other two.
% Exercise 25 gives angle 1 = angle 2, AC = BC and IJ = IK, so the plate is
% mirror-symmetric about CI: I is the MIDPOINT of AB (the draft had it at
% 0.46, which broke IJ = IK), and J, K sit symmetrically either side of it.
\begin{tikzpicture}[scale=1.479]
  \coordinate (A) at (0,0);
  \coordinate (B) at (3.22,0);
  \coordinate (C) at (1.61,1.92);
  \coordinate (I) at ($(A)!0.5!(B)$);
  \coordinate (J) at ($(A)!0.39!(B)$);
  \coordinate (K) at ($(A)!0.61!(B)$);
  \coordinate (D) at ($(A)!0.45!(C)$);
  \coordinate (E) at ($(B)!0.45!(C)$);
  \draw[given] (A) -- (B) -- (C) -- cycle;
  \draw[given] (C) -- (I);
  \draw[given] (A) -- (E);
  \draw[given] (B) -- (D);
  \draw[given] (D) -- (J);
  \draw[given] (E) -- (K);
  \coordinate (F) at (intersection of A--E and B--D);
  \coordinate (G) at (intersection of D--J and A--E);
  \coordinate (H) at (intersection of E--K and B--D);
  \node[vlab, left]  at (A) {$A$};
  \node[vlab, right] at (B) {$B$};
  \node[vlab, above] at (C) {$C$};
  \node[vlab, above left]  at (D) {$D$};
  \node[vlab, above right] at (E) {$E$};
  \node[vlab, below left] at (F) {$F$};
  % G and H sit just above the base between J/I and I/K; "below" laid their
  % letters straight across AB, so they go above their own points instead.
  \node[vlab] at ($(G)+(160:0.30)$) {$G$};
  \node[vlab] at ($(H)+(20:0.30)$) {$H$};
  \node[vlab, below] at (J) {$J$};
  \node[vlab, below] at (I) {$I$};
  \node[vlab, below] at (K) {$K$};
  \VIarc{B}{A}{E}{0.46cm}   \VInumb{B}{A}{E}{0.86cm}{1}
  \VIarc{A}{B}{D}{0.46cm}   \VInumb{A}{B}{D}{0.86cm}{2}
\end{tikzpicture}
\figcap{6--126}
\end{minipage}
\end{bkfigure}}

% ================================================================ Fig 6-127
% Book p.115, Algebra Review 4.  Triangle B C with apex A; D on AB, E on AC,
% and O the crossing of BE and CD.
\newcommand{\FIGVIONEHUNDREDTWENTYSEVEN}{\begin{bkfigure}[\linewidth]
\begin{tikzpicture}[scale=1.196]
  \coordinate (B) at (0,0);
  \coordinate (C) at (2.82,0);
  \coordinate (A) at (1.41,1.92);
  \coordinate (D) at ($(A)!0.46!(B)$);
  \coordinate (E) at ($(A)!0.46!(C)$);
  \draw[given] (B) -- (C) -- (A) -- cycle;
  \draw[given] (B) -- (E);
  \draw[given] (C) -- (D);
  \draw[given] (D) -- (E);
  \coordinate (O) at (intersection of B--E and C--D);
  \node[vlab, left]  at (B) {$B$};
  \node[vlab, right] at (C) {$C$};
  \node[vlab, above] at (A) {$A$};
  \node[vlab, above left]  at (D) {$D$};
  \node[vlab, above right] at (E) {$E$};
  \node[vlab, below] at (O) {$O$};
\end{tikzpicture}
\figcap{6--127}
\end{bkfigure}}
